%%%%%%%%%%%%%%%%%%%%%%%%%%%%%%%%%%%%%%%%%%%%%%%%%%%%%%%%%%%%%%%%%%
%%%%%%%% ICML 2016 EXAMPLE LATEX SUBMISSION FILE %%%%%%%%%%%%%%%%%
%%%%%%%%%%%%%%%%%%%%%%%%%%%%%%%%%%%%%%%%%%%%%%%%%%%%%%%%%%%%%%%%%%

% Use the following line _only_ if you're still using LaTeX 2.09.
%\documentstyle[icml2016,epsf,natbib]{article}
% If you rely on Latex2e packages, like most moden people use this:
\documentclass{article}

% use Times
\usepackage{times}
% For figures
\usepackage{graphicx} % more modern
%\usepackage{epsfig} % less modern
\usepackage{subfigure}

% For citations
\usepackage{natbib}

% For algorithms
% https://tex.stackexchange.com/questions/131646/algorithm2e-command-algorithm-already-defined
\usepackage{algorithm}
\usepackage{algorithmic}
\usepackage[algo2e,linesnumbered,lined,boxed,vlined]{algorithm2e}

\usepackage{amsfonts}
\usepackage{amssymb}


% As of 2011, we use the hyperref package to produce hyperlinks in the
% resulting PDF.  If this breaks your system, please commend out the
% following usepackage line and replace \usepackage{icml2016} with
% \usepackage[nohyperref]{icml2016} above.
\usepackage{hyperref}

% Packages hyperref and algorithmic misbehave sometimes.  We can fix
% this with the following command.
\newcommand{\theHalgorithm}{\arabic{algorithm}}

% Employ the following version of the ``usepackage'' statement for
% submitting the draft version of the paper for review.  This will set
% the note in the first column to ``Under review.  Do not distribute.''
%\usepackage{icml2016}
\usepackage[accepted]{icml2016}

% Employ this version of the ``usepackage'' statement after the paper has
% been accepted, when creating the final version.  This will set the
% note in the first column to ``Proceedings of the...''
%\usepackage[accepted]{icml2016}


%\usepackage{ifpdf}
%\usepackage{latexsym,amsfonts,amssymb,amsmath,graphicx,epsf,cite,bbm,float}
%\usepackage{ifpdf,flexisym}
%\usepackage{epstopdf}
%\usepackage{caption}
%\usepackage{subcaption}
\usepackage{comment}
\usepackage{amsthm}
\usepackage{mathtools}
%\usepackage{bbm}
%\usepackage{lettrine}
%\usepackage{algorithm}
%\usepackage{algorithmic}
%\usepackage[colorinlistoftodos,prependcaption]{todonotes}
\usepackage{tabularx}
\usepackage{multirow}

\newtheorem{theorem}{Theorem}
\newtheorem{lemma}{Lemma}
\newtheorem{proposition}{Proposition}
\newtheorem{corollary}{Corollary}
\newtheorem{definition}{Definition}
\newtheorem{example}{Example}
\newtheorem{remark}{Remark}
%\newtheorem{proof}{Proof}
\newtheorem{assumption}{Assumption}

% https://tex.stackexchange.com/questions/7793/defining-shortcuts
% I googled: overleaf define shortcut
\def\b{\ensuremath\boldsymbol}

% remove space between paragraphs:
% https://latex.org/forum/viewtopic.php?t=10877
\setlength{\parskip}{2pt}

% for adding page number:
% https://www.overleaf.com/learn/latex/Page_numbering
% https://tex.stackexchange.com/questions/151989/add-page-number-with-total-page-number-on-each-page
% https://tex.stackexchange.com/questions/56316/top-right-side-page-numbering
\usepackage{lastpage}
\usepackage{fancyhdr}
\fancyhead[R]{\thepage}

\usepackage{url}
\Urlmuskip=0mu plus 1mu  %--> for better url in references

% The \icmltitle you define below is probably too long as a header.
% Therefore, a short form for the running title is supplied here:
% \icmltitlerunning{Random Projection and the Johnson-Lindenstrauss Lemma: Tutorial and Survey}
\icmltitlerunning{}

\begin{document}

\AddToShipoutPictureBG*{%
  \AtPageUpperLeft{%
    \setlength\unitlength{1in}%
    \hspace*{\dimexpr0.5\paperwidth\relax}
    \makebox(0,-0.75)[c]{\normalsize {\color{black} To appear as a part of an upcoming textbook on dimensionality reduction and manifold learning.}}
    }}

\twocolumn[
\icmltitle{Johnson-Lindenstrauss Lemma, Linear and Nonlinear Random Projections, Random Fourier Features, and Random Kitchen Sinks: Tutorial and Survey}

% It is OKAY to include author information, even for blind
% submissions: the style file will automatically remove it for you
% unless you've provided the [accepted] option to the icml2016
% package.
\icmlauthor{Benyamin Ghojogh}{bghojogh@uwaterloo.ca}
\icmladdress{Department of Electrical and Computer Engineering, 
\\Machine Learning Laboratory, University of Waterloo, Waterloo, ON, Canada}
\icmlauthor{Ali Ghodsi}{ali.ghodsi@uwaterloo.ca}
\icmladdress{Department of Statistics and Actuarial Science \& David R. Cheriton School of Computer Science, 
\\Data Analytics Laboratory, University of Waterloo, Waterloo, ON, Canada}
\icmlauthor{Fakhri Karray}{karray@uwaterloo.ca}
\icmladdress{Department of Electrical and Computer Engineering, 
\\Centre for Pattern Analysis and Machine Intelligence, University of Waterloo, Waterloo, ON, Canada}
\icmlauthor{Mark Crowley}{mcrowley@uwaterloo.ca}
\icmladdress{Department of Electrical and Computer Engineering, 
\\Machine Learning Laboratory, University of Waterloo, Waterloo, ON, Canada}


% You may provide any keywords that you
% find helpful for describing your paper; these are used to populate
% the "keywords" metadata in the PDF but will not be shown in the document
\icmlkeywords{Tutorial}

\vskip 0.3in
]

\begin{abstract}
This is a tutorial and survey paper on the Johnson-Lindenstrauss (JL) lemma and linear and nonlinear random projections. We start with linear random projection and then justify its correctness by JL lemma and its proof. Then, sparse random projections with $\ell_1$ norm and interpolation norm are introduced. Two main applications of random projection, which are low-rank matrix approximation and approximate nearest neighbor search by random projection onto hypercube, are explained. Random Fourier Features (RFF) and Random Kitchen Sinks (RKS) are explained as methods for nonlinear random projection. Some other methods for nonlinear random projection, including extreme learning machine, randomly weighted neural networks, and ensemble of random projections, are also introduced. 
\end{abstract}



\section{Introduction}

Linear dimensionality reduction methods project data onto the low-dimensional column space of a projection matrix.
Many of these methods, such as Principal Component Analysis (PCA) \cite{ghojogh2019unsupervised} and Fisher Discriminant Analysis (FDA) \cite{ghojogh2019fisher}, learn a projection matrix for either better representation of data or discrimination of classes in the subspace. 
For example, PCA learns the projection matrix to maximize the variance of projected data. FDA learns the projections matrix to discriminate classes in the subspace. 
However, it was seen that if we do not learn the projection matrix and just sample the elements of projection matrix randomly, it still works. 
In fact, random projection preserves the distances of points with a small error after projection. 
Hence, this random projection works very well surprisingly although its projection matrix is not learned! 
For justifying why random projection works well, the Johnson-Lindenstrauss (JL) lemma \cite{johnson1984extensions} was proposed which bounds the error of random projection. 
Random projection is a probabilistic dimensionality reduction method. 
Theories for various random projection methods have a similar approach (e.g., see the proofs in this paper) and similar bounds. 
In terms of dealing with probability bounds, it can be slightly related to the Probably Approximately Correct (PAC) learning \cite{shalev2014understanding}. 
A survey on linear random projection is \cite{xie2017survey} and a good book on this topic is \cite{vempala2005random}. Slides \cite{liberty2006random,ward2014dimension} are also good slides about linear random projection. 

The theory of linear random projection has been developed during years; however, the theory of nonlinear random projection still needs to be improved because it is more complicated to analyze than linear random projection. 
Nonlinear random projection can be modeled as a linear random projection followed by a nonlinear function. 
Random Fourier Features (RFF) \cite{rahimi2007random} and Random Kitchen Sinks (RKS) \cite{rahimi2008weighted,rahimi2008uniform} are two initial works on nonlinear random projection. Some other methods have also been proposed for nonlinear random projection. They are Extreme Learning Machine (ELM) \cite{huang2006extreme,huang2011extreme,tang2015extreme}, randomly weighted neural network \cite{jarrett2009best}, and ensemble random projections \cite{schclar2009random,cannings2015random,karimi2017ensembles}. 

This is a tutorial and survey paper on linear and nonlinear random projection.
The remainder of this paper is organized as follows. We introduce linear projection and linear random projection in Section \ref{section_linear_random_projection}. The JL lemma and its proof are provided in Section \ref{section_JL_lemma}. Then, sparse random projection by $\ell_1$ norm is introduced in Section \ref{section_sparse_linear_random_projection}. Applications of linear random projection, including low-rank matrix approximation and approximate nearest neighbor search, are explained in Section \ref{section_applications_of_random_projection}. Section \ref{section_RFF_and_RKS} introduce the theory of RFF and RKS for nonlinear random projection. Some other methods for nonlinear random projection, including ELM, random neural network, and ensemble of random projections, are introduced in Section \ref{section_other_nonlinear_random_projections}. Finally, Section \ref{section_conclusion} concludes the paper. 

\section*{Required Background for the Reader}

This paper assumes that the reader has general knowledge of calculus, probability, linear algebra, and basics of optimization. 

\section{Linear Random Projection}\label{section_linear_random_projection}

\subsection{Linear Projection}

\subsubsection{Projection Formulation}

Assume we have a data point $\b{x} \in \mathbb{R}^d$. We want to project this data point onto the vector space spanned by $p$ vectors $\{\b{u}_1, \dots, \b{u}_p\}$ where each vector is $d$-dimensional and usually $p \ll d$. We stack these vectors column-wise in matrix $\b{U} = [\b{u}_1, \dots, \b{u}_p] \in \mathbb{R}^{d \times p}$. In other words, we want to project $\b{x}$ onto the column space of $\b{U}$, denoted by $\mathbb{C}\text{ol}(\b{U})$.

The projection of $\b{x} \in \mathbb{R}^d$ onto $\mathbb{C}\text{ol}(\b{U}) \in \mathbb{R}^p$ and then its representation in the $\mathbb{R}^d$ (its reconstruction) can be seen as a linear system of equations:
\begin{align}\label{equation_projection}
\mathbb{R}^d \ni \widehat{\b{x}} := \b{U \beta},
\end{align}
where we should find the unknown coefficients $\b{\beta} \in \mathbb{R}^p$. 

If the $\b{x}$ lies in the $\mathbb{C}\text{ol}(\b{U})$ or $\textbf{span}\{\b{u}_1, \dots, \b{u}_p\}$, this linear system has exact solution, so $\widehat{\b{x}} = \b{x} = \b{U \beta}$. However, if $\b{x}$ does not lie in this space, there is no any solution $\b{\beta}$ for $\b{x} = \b{U \beta}$ and we should solve for projection of $\b{x}$ onto $\mathbb{C}\text{ol}(\b{U})$ or $\textbf{span}\{\b{u}_1, \dots, \b{u}_p\}$ and then its reconstruction. In other words, we should solve for Eq. (\ref{equation_projection}). In this case, $\widehat{\b{x}}$ and $\b{x}$ are different and we have a residual:
\begin{align}\label{equation_residual_1}
\b{r} = \b{x} - \widehat{\b{x}} = \b{x} - \b{U \beta},
\end{align}
which we want to be small. As can be seen in Fig. \ref{figure_residual_and_space}, the smallest residual vector is orthogonal to $\mathbb{C}\text{ol}(\b{U})$; therefore:
\begin{align}
\b{x} - \b{U\beta} \perp \b{U} &\implies \b{U}^\top (\b{x} - \b{U \beta}) = 0, \nonumber \\
& \implies \b{\beta} = (\b{U}^\top \b{U})^{-1} \b{U}^\top \b{x}. \label{equation_beta}
\end{align}
It is noteworthy that the Eq. (\ref{equation_beta}) is also the formula of coefficients in linear regression where the input data are the rows of $\b{U}$ and the labels are $\b{x}$; however, our goal here is different. 

\begin{figure}[!t]
\centering
\includegraphics[width=2.2in]{./images/residual_and_space}
\caption{The residual and projection onto the column space of $\b{U}$. Credit of image is for \cite{ghojogh2019unsupervised}.}
\label{figure_residual_and_space}
\end{figure}

Plugging Eq. (\ref{equation_beta}) in Eq. (\ref{equation_projection}) gives us:
\begin{align*}
\widehat{\b{x}} = \b{U} (\b{U}^\top \b{U})^{-1} \b{U}^\top \b{x}.
\end{align*}
We define:
\begin{align}\label{equation_hat_matrix}
\mathbb{R}^{d \times d} \ni \b{\Pi} := \b{U} (\b{U}^\top \b{U})^{-1} \b{U}^\top,
\end{align}
as ``projection matrix'' because it projects $\b{x}$ onto $\mathbb{C}\text{ol}(\b{U})$ (and reconstructs back).
Note that $\b{\Pi}$ is also referred to as the ``hat matrix'' in the literature because it puts a hat on top of $\b{x}$.

If the vectors $\{\b{u}_1, \dots, \b{u}_p\}$ are orthonormal (the matrix $\b{U}$ is orthogonal), we have $\b{U}^\top = \b{U}^{-1}$ and thus $\b{U}^\top \b{U} = \b{I}$. Therefore, Eq. (\ref{equation_hat_matrix}) is simplified:
\begin{align}
& \b{\Pi} = \b{U} \b{U}^\top.
\end{align}
So, we have:
\begin{align}\label{equation_x_hat}
\widehat{\b{x}} = \b{\Pi}\, \b{x} = \b{U} \b{U}^\top \b{x}.
\end{align}

\subsubsection{Projection onto a Subspace}\label{section_projection_subspace}

In subspace learning, the projection of a vector $\b{x} \in \mathbb{R}^d$ onto the column space of $\b{U} \in \mathbb{R}^{d \times p}$ (a $p$-dimensional subspace spanned by $\{\b{u}_j\}_{j=1}^p$ where $\b{u}_j \in \mathbb{R}^d$) is defined as:
\begin{align}
&\mathbb{R}^{p} \ni \widetilde{\b{x}} := \b{U}^\top \b{x}, \label{equation_projection_training_onePoint_severalDirections} \\
&\mathbb{R}^{d} \ni \widehat{\b{x}} := \b{U}\b{U}^\top \b{x} = \b{U} \widetilde{\b{x}}, \label{equation_reconstruction_training_onePoint_severalDirections}
\end{align}
where $\widetilde{\b{x}}$ and $\widehat{\b{x}}$ denote the projection and reconstruction of $\b{x}$, respectively.

If we have $n$ data points, $\{\b{x}_i\}_{i=1}^n$, which can be stored column-wise in a matrix $\b{X} \in \mathbb{R}^{d \times n}$, the projection and reconstruction of $\b{X}$ are defined as:
\begin{align}
&\mathbb{R}^{p \times n} \ni \widetilde{\b{X}} := \b{U}^\top \b{X}, \label{equation_projection_training_SeveralPoints_severalDirections} \\
&\mathbb{R}^{d \times n} \ni \widehat{\b{X}} := \b{U}\b{U}^\top \b{X} = \b{U} \widetilde{\b{X}}, \label{equation_reconstruction_training_SeveralPoints_severalDirections}
\end{align}
respectively.

If we have an out-of-sample data point $\b{x}_t$ which was not used in calculation of $\b{U}$, the projection and reconstruction of it are defined as:
\begin{align}
&\mathbb{R}^{p} \ni \widetilde{\b{x}}_t := \b{U}^\top \b{x}_t, \label{equation_projection_outOfSample_onePoint_severalDirections} \\
&\mathbb{R}^{d} \ni \widehat{\b{x}}_t := \b{U}\b{U}^\top \b{x}_t = \b{U} \widetilde{\b{x}}_t, \label{equation_reconstruction_outOfSample_onePoint_severalDirections}
\end{align}
respectively.

In case we have $n_t$ out-of-sample data points, $\{\b{x}_{t,i}\}_{i=1}^{n_t}$, which can be stored column-wise in a matrix $\b{X}_t \in \mathbb{R}^{d \times n_t}$, the projection and reconstruction of $\b{X}_t$ are defined as:
\begin{align}
&\mathbb{R}^{p \times n_t} \ni \widetilde{\b{X}}_t := \b{U}^\top \b{X}_t, \label{equation_projection_outOfSample_SeveralPoints_severalDirections} \\
&\mathbb{R}^{d \times n_t} \ni \widehat{\b{X}}_t := \b{U}\b{U}^\top \b{X}_t = \b{U} \widetilde{\b{X}}_t, \label{equation_reconstruction_outOfSample_SeveralPoints_severalDirections}
\end{align}
respectively.

For the properties of the projection matrix $\b{U}$, refer to \cite{ghojogh2019unsupervised}.

\subsection{Linear Random Projection}

Linear random projection is projection of data points onto the column space of a projection matrix where the elements of projection matrix are i.i.d. random variables sampled from a distribution with zero mean and (possibly scaled) unit variance. In other words, random projection is a function $f\!: \mathbb{R}^d \rightarrow \mathbb{R}^p$, $f\!: \b{x} \mapsto \b{U}^\top \b{x}$:
\begin{align}\label{equation_linear_random_projection}
\mathbb{R}^p \ni f(\b{x}) := \b{U}^\top \b{x} = \sum_{t=1}^p \b{u}_t^\top \b{x} = \sum_{j=1}^d \sum_{t=1}^p u_{jt}\, x_j,
\end{align}
where $\b{U} = [\b{u}_1, \dots, \b{u}_p] \in \mathbb{R}^{d \times p}$ is the random projection matrix, $u_{jt}$ is the $(j,t)$-th element of $\b{U}$, and $x_j$ is the $j$-th element of $\b{x} \in \mathbb{R}^d$. 
The elements of $\b{U}$ are sampled from a distribution with zero mean and (possibly scaled) unit variance. For example, we can use the Gaussian distribution $u_{j,t} \sim \mathcal{N}(0, 1), \forall j,t$. 
Some example distributions to use for random projection are Gaussian \cite{giryes2016deep} and Cauchy \cite{li2007nonlinear,ramirez2012reconstruction} distributions. 

It is noteworthy that, in some papers, the projection is normalized:
\begin{align}\label{equation_linear_random_projection_normalized}
f\!: \b{x} \mapsto \frac{1}{\sqrt{p}} \b{U}^\top \b{x}.
\end{align}
If we have a dataset of $d$-dimensional points with sample size $n$, we can stack the points in $\b{X} := [\b{x}_1, \dots, \b{x}_n] \in \mathbb{R}^{d \times n}$. 
Eq. (\ref{equation_linear_random_projection}) or (\ref{equation_linear_random_projection_normalized}) is stated as:
\begin{align}
\mathbb{R}^{p \times n} \ni f(\b{X}) := \b{U}^\top \b{X} \quad \text{or} \quad \frac{1}{\sqrt{p}} \b{U}^\top \b{X}.
\end{align}

In contrast to other linear dimensionality reduction methods which learn the projection matrix using training dataset for better data representation of class discrimination, random projection does not learn the projection matrix but randomly samples it independent of data.  
Surprisingly, projection of data onto the column space of this random matrix works very well although the projection matrix is completely random and independent of data. 
The JL lemma \cite{johnson1984extensions} justifies why random projection works. In the following, we introduce this lemma and its proof. 

\section{The Johnson-Lindenstrauss Lemma}\label{section_JL_lemma}

% https://en.wikipedia.org/wiki/Johnson%E2%80%93Lindenstrauss_lemma

\begin{theorem}[Johnson-Lindenstrauss Lemma \cite{johnson1984extensions}]\label{theorem_JL_lemma}
For any set $\mathcal{X} := \{\b{x}_i \in \mathbb{R}^d\}_{i=1}^n$, any integer $n$ as the sample size, and any $0<\epsilon<1$ as error tolerance, let $p$ be a positive integer satisfying:
\begin{align}\label{equation_JL_lemma_p_lower_bound}
p \geq \Omega(\frac{\ln(n)}{\epsilon^2 - \epsilon^3}),
\end{align}
where $\ln(.)$ is the natural logarithm and $\Omega(.)$ is the lower bound complexity [n.b. some works state Eq. (\ref{equation_JL_lemma_p_lower_bound}) as:
\begin{align}\label{equation_JL_lemma_p_lower_bound_relaxed}
p \geq \Omega(\epsilon^{-2} \ln(n)).
\end{align}
by ignoring $\epsilon^3$ against $\epsilon^2$ in the denominator because $\epsilon \in (0,1)$]. 
There exists a linear map $f\!: \mathbb{R}^d \rightarrow \mathbb{R}^p$, $f\!: \b{x} \mapsto \b{U}^\top \b{x}$, with projection matrix $\b{U} = [\b{u}_1, \dots, \b{u}_p] \in \mathbb{R}^{d \times p}$, such that we have:
\begin{equation}\label{equation_JL_lemma_correct_projection}
\begin{aligned}
(1 - \epsilon) \|\b{x}_i - \b{x}_j\|_2^2 \leq \|&f(\b{x}_i) - f(\b{x}_j)\|_2^2 \\
&\leq (1 + \epsilon) \|\b{x}_i - \b{x}_j\|_2^2,
\end{aligned}
\end{equation}
for all $\b{x}_i, \b{x}_j \in \mathcal{X}$, with probability of success as:
\begin{align}\label{equation_random_projection_prob_success}
\mathbb{P}\Big(
(1 - \epsilon) \|\b{x}_i - &\b{x}_j\|_2^2 \leq \|f(\b{x}_i) - f(\b{x}_j)\|_2^2 \nonumber \\
&\leq (1 + \epsilon) \|\b{x}_i - \b{x}_j\|_2^2 
\Big) \geq 1 - \delta,
\end{align}
where $\delta := 2 e^{-(\epsilon^2 - \epsilon^3) (p/4)}$ and the elements of the projection matrix are i.i.d. random variables with mean zero and (scaled) unit variance. An example is $u_{ij} \sim \mathcal{N}(0, 1/p) = (1/\sqrt{p}) \mathcal{N}(0, 1)$ where $u_{ij}$ denotes the $(i,j)$-th element of $\b{U}$. 
\end{theorem}
\begin{proof}
The proof is based on \cite{karimi2018exploring,indyk1998approximate,dasgupta1999elementary,dasgupta2003elementary,achlioptas2003database,achlioptas2001database,matouvsek2008variants}.
Also some other proofs exists for JL lemma, such as \cite{frankl1988johnson,arriaga1999algorithmic}. 

\hfill \break
The proof can be divided into three steps:

\textbf{Step 1 of proof:} 
In this step, we show that the random projection preserves local distances of points in expectation.
Let $\b{x}_i = [x_{i1}, \dots, x_{id}]^\top \in \mathbb{R}^d$. We have:
\begin{align}
&\mathbb{E}\big[\|f(\b{x}_i) - f(\b{x}_j)\|_2^2\big] = \mathbb{E}\big[\|\b{U}^\top \b{x}_i - \b{U}^\top \b{x}_j\|_2^2\big] \nonumber\\
&\overset{(a)}{=} \mathbb{E}\Bigg[\sum_{t=1}^p \big(\b{u}_t^\top \b{x}_i - \b{u}_t^\top \b{x}_j\big)^2\Bigg] \nonumber\\
&\overset{(b)}{=} \sum_{t=1}^p \mathbb{E}\big[\big(\b{u}_t^\top \b{x}_i - \b{u}_t^\top \b{x}_j\big)^2\big] \nonumber\\
&= \sum_{t=1}^p \Big( \mathbb{E}\big[\big(\b{u}_t^\top \b{x}_i)^2\big] + \mathbb{E}\big[(\b{u}_t^\top \b{x}_j\big)^2\big] \nonumber\\
&~~~~~~~~~~~ - 2\,\mathbb{E}\big[(\b{u}_t^\top \b{x}_i \b{u}_t^\top \b{x}_j\big)^2\big] \Big), \label{equation_expectation_difference_of_mapping}
\end{align}
where $(a)$ is because of definition of $\ell_2$ norm and $(b)$ is because expectation is a linear operator. 
We know $u_{ij}$ has zero mean and unit variance. Let its unit variance be normalized as $1/p$. Then, the covariance of elements of projection matrix is:
\begin{align}\label{equation_covariance_of_projection_matrix}
\mathbb{E}[u_{kt} u_{lt}] = \delta_{kl} := 
\left\{
    \begin{array}{ll}
        0 & \mbox{if } i \neq j \\
        1/p & \mbox{if } i = j,
    \end{array}
\right.
\end{align}
which is a normalized Kronekcer delta.
The first term in this expression is:
\begin{align*}
&\mathbb{E}\big[(\b{u}_t^\top \b{x}_i)^2\big] = \mathbb{E}\big[\big(\sum_{k=1}^d u_{kt} x_{ik}\big)^2\big] \\
&= \mathbb{E}\Big[\sum_{k=1}^d \sum_{l=1}^d u_{kt} u_{lt} x_{ik} x_{il}\Big] \overset{(a)}{=} \sum_{k=1}^d \sum_{l=1}^d \mathbb{E}[u_{kt} u_{lt} x_{ik} x_{il}] \\
&\overset{(b)}{=} \sum_{k=1}^d \sum_{l=1}^d x_{ik} x_{il} \mathbb{E}[u_{kt} u_{lt}] \overset{(\ref{equation_covariance_of_projection_matrix})}{=} \sum_{k=1}^d \sum_{l=1}^d x_{ik} x_{il} \delta_{kl} \\
&= \frac{1}{p} \sum_{k=1}^d x_{ik}^2 = \frac{\|\b{x}_i\|_2^2}{p},
\end{align*}
where $(a)$ is because of linearity of expectation and $(b)$ is because data are deterministic.
Likewise, the second term of Eq. (\ref{equation_expectation_difference_of_mapping}) is $\mathbb{E}\big[(\b{u}_t^\top \b{x}_j)^2\big] = \|\b{x}_j\|_2^2 / p$. 
The third term is:
\begin{align*}
&\mathbb{E}\big[(\b{u}_t^\top \b{x}_i \b{u}_t^\top \b{x}_j\big)^2\big] = \mathbb{E}\big[\big(\sum_{k=1}^d u_{kt} x_{ik}\big) \big(\sum_{l=1}^d u_{lt} x_{jl}\big)\big] \\
&= \mathbb{E}\Big[\sum_{k=1}^d \sum_{l=1}^d u_{kt} u_{lt} x_{ik} x_{jl}\Big] = \sum_{k=1}^d \sum_{l=1}^d \mathbb{E}[u_{kt} u_{lt} x_{ik} x_{jl}] \\
&= \sum_{k=1}^d \sum_{l=1}^d x_{ik} x_{jl} \mathbb{E}[u_{kt} u_{lt}] \overset{(\ref{equation_covariance_of_projection_matrix})}{=} \sum_{k=1}^d \sum_{l=1}^d x_{ik} x_{jl} \delta_{kl} \\
&= \frac{1}{p} \sum_{k=1}^d x_{ik} x_{jk} = \frac{\b{x}_i^\top \b{x}_j}{p}.
\end{align*}
Therefore, Eq. (\ref{equation_expectation_difference_of_mapping}) is simplified to:
\begin{align}
\mathbb{E}\big[\|f(\b{x}_i) &- f(\b{x}_j)\|_2^2\big] \nonumber\\
&= \sum_{t=1}^p \Big( \frac{\|\b{x}_i\|_2^2}{p} + \frac{\|\b{x}_j\|_2^2}{p} - 2\frac{\b{x}_i^\top \b{x}_j}{p} \Big) \nonumber\\
&= \sum_{t=1}^p \Big( \frac{\|\b{x}_i - \b{x}_j\|_2^2}{p} \Big) = \|\b{x}_i - \b{x}_j\|_2^2 \sum_{t=1}^p \frac{1}{p} \nonumber\\
&= \|\b{x}_i - \b{x}_j\|_2^2, \quad \forall i, j \in \{1, \dots, n\}.
\end{align}
This shows that this projection by random projection matrix preserves the local distances of data points in expectation. 

\hfill \break
\textbf{Step 2 of proof:}
In this step, we show that the variance of local distance preservation is bounded and small for one pair of points. The probability of error in local distance preservation is:
\begin{align*}
&\mathbb{P}\Big(\|f(\b{x}_i) - f(\b{x}_j)\|_2^2 > (1+\epsilon) \|\b{x}_i - \b{x}_j\|_2^2\Big) \\
&= \mathbb{P}\Big(\|\b{U}^\top \b{x}_i - \b{U}^\top \b{x}_j\|_2^2 > (1+\epsilon) \|\b{x}_i - \b{x}_j\|_2^2\Big) \\
&= \mathbb{P}\Big(\|\b{U}^\top (\b{x}_i - \b{x}_j)\|_2^2 > (1+\epsilon) \|\b{x}_i - \b{x}_j\|_2^2\Big).
\end{align*}
Let $\b{x}_D := \b{x}_i - \b{x}_j$ so:
\begin{align*}
\mathbb{P}\Big(\|\b{U}^\top \b{x}_D\|_2^2 > (1+\epsilon) \|\b{x}_d\|_2^2\Big).
\end{align*}
We know that $u_{ij} \sim (1/\sqrt{p}) \mathcal{N}(0, 1)$. 
So, the projected data are $\mathbb{R}^p \ni \b{y}_D = \b{U}^\top \b{x}_D$ and we have $\b{y}_D \sim (1/\sqrt{p}) \mathcal{N}(0, \|\b{x}_D\|_2^2) = \mathcal{N}(0, (1/p)\|\b{x}_D\|_2^2)$ because of the quadratic characteristic of variance. 

We define a matrix $\b{A} = [a_{ij}] \in \mathbb{R}^{d \times p}$ where $a_{ij} \sim \mathcal{N}(0, 1)$. Assume $\b{z}_D = [z_1, \dots, z_p]^\top := (\b{A}^\top \b{x}_D) / \|\b{x}_D\|_2$. Hence:
\begin{align*}
&\mathbb{P}\Big(\|\frac{1}{\sqrt{p}}\b{A}^\top \b{x}_D\|_2^2 > (1+\epsilon) \|\b{x}_D\|_2^2\Big) \\
&= \mathbb{P}\Big(\frac{1}{p}\|\b{A}^\top \b{x}_D\|_2^2 > (1+\epsilon) \|\b{x}_D\|_2^2\Big) \\
&\overset{(a)}{=} \mathbb{P}\Big(\frac{1}{\|\b{x}_D\|_2^2} \|\b{A}^\top \b{x}_D\|_2^2 > (1+\epsilon) p\Big) \\
&= \mathbb{P}\Big(\big\|\frac{\b{A}^\top \b{x}_D}{\|\b{x}_D\|_2}\big\|_2^2 > (1+\epsilon) p\Big) \\
&\overset{(b)}{=} \mathbb{P}\Big(\|\b{z}_D\|_2^2 > (1+\epsilon) p\Big) \\
&= \mathbb{P}\Big(\sum_{k=1}^p z_k^2 > (1+\epsilon)\, p\Big) \overset{(c)}{=} \mathbb{P}\Big(\chi_p^2 > (1+\epsilon)\, p\Big),
\end{align*}
where $(a)$ is because the sides are multiplied by $p / \|\b{x}_D\|_2^2$, $(b)$ is because of the definition of $\b{z}_D$, and $(c)$ is because the summation of $p$ squared standard normal distributions is the chi-squared distribution with $p$ degrees of freedom, denoted by $\chi_p^2$. 
We have:
\begin{align}
\mathbb{P}\Big(\chi_p^2 > &(1+\epsilon)\, p\Big) \overset{(a)}{=} \mathbb{P}\Big(e^{\lambda \chi_p^2} > e^{\lambda (1+\epsilon)\, p}\Big) \nonumber\\
&\overset{(b)}{\leq} \frac{\mathbb{E}[e^{\lambda \chi_p^2}]}{e^{\lambda (1+\epsilon)\, p}} \overset{(c)}{=} \frac{(1-2\lambda)^{-p/2}}{e^{\lambda (1+\epsilon)\, p}} \label{equation_prob_error_1}
\end{align}
where $(a)$ is because of taking power of $e$ with parameter $\lambda \geq 0$, $(b)$ is because of the Markov inequality (i.e., $\mathbb{P}(x \geq a) \leq \mathbb{E}[x]/a$), and $(c)$ is because of the moment generating function of the chi-squared distribution $x \sim \chi_p^2$ which is $M_x(t) = \mathbb{E}[e^{\lambda x}] = (1-2\lambda)^{-p/2}$. 
We want to minimize the probability of error. So, we minimize this error in Eq. (\ref{equation_prob_error_1}) with respect to the parameter $\lambda$:
\begin{align*}
\frac{\partial}{\partial \lambda} (\frac{(1-2\lambda)^{-p/2}}{e^{\lambda (1+\epsilon)\, p}}) \overset{\text{set}}{=} 0 \implies \lambda = \frac{\epsilon}{2(1+\epsilon)}.
\end{align*}
Substituting this equation in Eq. (\ref{equation_prob_error_1}) gives:
\begin{align*}
\mathbb{P}&\Big(\chi_p^2 > (1+\epsilon)\, p\Big) = \frac{(1+\epsilon)^{p/2}}{e^{(p \epsilon) / 2}} \\
&= ((1+\epsilon) e^{-\epsilon})^{p/2} \overset{(a)}{\leq} e^{-(p/4) (\epsilon^2 - \epsilon^3)},
\end{align*}
where $(a)$ is because $1+\epsilon \leq e^{(\epsilon - (\epsilon^2 - \epsilon^3)/2)}$. 
Above, we derived the probability of $\mathbb{P}(\|f(\b{x}_i) - f(\b{x}_j)\|_2^2 > (1+\epsilon) \|\b{x}_i - \b{x}_j\|_2^2)$.
Likewise, we can derive the same amount of probability for $\mathbb{P}(\|f(\b{x}_i) - f(\b{x}_j)\|_2^2 < (1-\epsilon) \|\b{x}_i - \b{x}_j\|_2^2)$. 
For the projection of pair $\b{x}_i$ and $\b{x}_j$, we want to have:
\begin{align}
\|f(\b{x}_i) - f(\b{x}_j)\|_2^2 \approx \|\b{x}_i - \b{x}_j\|_2^2.
\end{align}
Hence, the total probability of error for the projection of pair $\b{x}_i$ and $\b{x}_j$ is:
\begin{align}
\mathbb{P}\Big(&\Big[\|f(\b{x}_i) - f(\b{x}_j)\|_2^2 > (1+\epsilon) \|\b{x}_i - \b{x}_j\|_2^2\Big] \nonumber\\
&\cup
\Big[\|f(\b{x}_i) - f(\b{x}_j)\|_2^2 < (1-\epsilon) \|\b{x}_i - \b{x}_j\|_2^2\Big] \Big) \nonumber\\
&\leq 2\, e^{-(p/4) (\epsilon^2 - \epsilon^3)}. \label{equation_JL_prob_error_twoPoints}
\end{align}
This shows that the probability of error for every pair $\b{x}_i$ and $\b{x}_j$ is bounded. This also proves the Eq. (\ref{equation_random_projection_prob_success}).

\hfill \break
\textbf{Step 3 of proof:}
In this step, we bound the probability of error for all pairs of data points:
We use the Bonferroni's inequality or the so-called union bound \cite{bonferroni1936teoria}:
\begin{align}\label{equation_Bonferroni_inequality}
\mathbb{P}\Big(\bigcup_i E_i\Big) \leq \sum_{i} \mathbb{P}(E_i),
\end{align}
where $E_i$ is the $i$-th probability event. We have $\binom{n}{2} = n(n-1)/2$ pairs of points. Hence, the total probability of error for all pairs of points is:
\begin{align}
\mathbb{P}(\text{error}) &\leq \frac{n(n-1)}{2} \times 2\, e^{-(p/4) (\epsilon^2 - \epsilon^3)} \nonumber\\
&= n(n-1) e^{-(p/4) (\epsilon^2 - \epsilon^3)} \leq \delta', \label{equation_prob_error_JL_lemma}
\end{align}
where the upper-bound $0 \leq \delta' \ll 1$ can be selected according to $n$, $p$, and $\epsilon$. Therefore, the probability of error is bounded. This shows that Eq. (\ref{equation_JL_lemma_correct_projection}) holds with probability $1-\delta'$. Q.E.D.
\end{proof}

It is noteworthy that Eq. (\ref{equation_random_projection_prob_success}) can be restated as:
\begin{align}
\mathbb{P}\Big(
\Big| \|f(\b{x}_i) - f(\b{x}_j)\|_2^2 - \|\b{x}_i - &\b{x}_j\|_2^2 \Big| \geq \epsilon \Big) \leq \delta.
\end{align}
This shows that there is an upper-bound on the probability of error in random projection. Therefore, after linear random projection, the distance is preserved after projection for every pair of points. 

The following lemma proves the required lower-bound on the dimensionality of subspace for random projection to work well. 
% slides \cite{liberty2006random}, slides 5 to 12, especially slide 12:
\begin{lemma}[Lower bound on dimensionality of subspace \cite{indyk1998approximate}]
The dimensionality of subspace $p$ should satisfy Eq. (\ref{equation_JL_lemma_p_lower_bound}).
\end{lemma}
\begin{proof}
The probability of error, which is Eq. (\ref{equation_prob_error_JL_lemma}), should be smaller than a constant $\delta$:
\begin{align*}
&\mathbb{P}(\text{error}) = n(n-1) e^{-(p/4) (\epsilon^2 - \epsilon^3)} \leq \delta' \\
&\implies -(p/4) (\epsilon^2 - \epsilon^3) \leq \ln(\frac{\delta'}{n(n-1)}) \\
&\implies (p/4) (\epsilon^2 - \epsilon^3) \geq \ln(\frac{n(n-1)}{\delta'}) \\
&\implies (p/4) \geq \frac{1}{\epsilon^2 - \epsilon^3} \ln(\frac{n(n-1)}{\delta'}) \approx \frac{\ln(\frac{n^2}{\delta'})}{\epsilon^2 - \epsilon^3} \\
&\implies p \geq \frac{8\ln(n) - 4 \ln(\delta')}{\epsilon^2 - \epsilon^3} = \Omega(\frac{\ln(n)}{\epsilon^2 - \epsilon^3}).
\end{align*}
Q.E.D.
\end{proof}

% https://en.wikipedia.org/wiki/Concentration_of_measure
% \cite{ward2014dimension}, slide 12
\begin{definition}[Concentration of Measure \cite{talagrand1996new}]\label{definition_concentration_of_measure}
Inspired by Eq. (\ref{equation_random_projection_prob_success}), we have the inequality {\citep[Lemma A.1]{plan2013one}}:
\begin{align}
\mathbb{P}\Big((1-\epsilon) \|\b{x}\|_2^2 \leq \|\b{U}^\top \b{x}\|_2^2 \leq& (1+\epsilon) \|\b{x}\|_2^2\Big) \nonumber\\
&\geq 1 - 2e^{-c \epsilon^2 p}, \label{equation_concentration_of_measure}
\end{align}
where $c>0$ is a constant, $\epsilon$ is the error tolerance, and $p$ is the dimensionality of projected data by the projection matrix $\b{U} \in \mathbb{R}^{d \times p}$.
This is the concentration of measure for a Gaussian random matrix. This inequality holds for random projection.
\end{definition}
Comparing Eqs. (\ref{equation_random_projection_prob_success}) and (\ref{equation_concentration_of_measure}) shows that random projection is equivalent to concentration of measure.

In the following, we provide a lemma which is going to be used in the proof of Lemma \ref{lemma_prob_success_whp}. Lemma \ref{lemma_prob_success_whp} provides the order of probability of success for random projection and we will use it to explain why random projection is correct (i.e., works well) with high probability. 
\begin{lemma}[\cite{dasgupta1999elementary,dasgupta2003elementary}]\label{lemma_P_beta_and_exp}
Let $p<d$ and $L := \|\b{z}_D\|_2^2$ where $\b{z}_D$ is the random projection of $\b{x}_D := \b{x}_i - \b{x}_j$ onto a $p$-dimensional subspace. If $\beta < 1$, then:
\begin{align}\label{equation_P_beta_and_exp_1}
\mathbb{P}(L \leq \beta p/d) \leq \exp\big(\frac{p}{2} (1 - \beta + \ln \beta)\big).
\end{align}
If $\beta > 1$, then:
\begin{align}\label{equation_P_beta_and_exp_2}
\mathbb{P}(L \geq \beta p/d) \leq \exp\big(\frac{p}{2} (1 - \beta + \ln \beta)\big).
\end{align}
\end{lemma}
\begin{proof}
Refer to \cite{dasgupta1999elementary,dasgupta2003elementary}.
\end{proof}

\begin{lemma}[\cite{dasgupta1999elementary,dasgupta2003elementary}]\label{lemma_prob_success_whp}
If $p$ satisfies Eq. (\ref{equation_JL_lemma_p_lower_bound}) as:
\begin{align}\label{equation_JL_lemma_p_lower_bound_2}
p \geq 4(\frac{\epsilon^2}{2} - \frac{\epsilon^3}{3})^{-1} \ln(n) = \Omega(\frac{\ln(n)}{\epsilon^2 - \epsilon^3}),
\end{align}
the local distances of points are not distorted by random projection no more than $(1 \pm \epsilon)$ with probability $\mathcal{O}(1/n^2)$. In other words, the probability of error in random projection is of the order $\mathcal{O}(1/n^2)$.
\end{lemma}
\begin{proof}
The proof is based on \cite{dasgupta2003elementary}.
\hfill\break
\textbf{Step 1 of proof:}
Consider one side of error which is $\mathbb{P}(\|f(\b{x}_i) - f(\b{x}_j)\|_2^2 \leq (1-\epsilon) (p/d) \|\b{x}_i - \b{x}_j\|_2^2)$. Take $L = \|f(\b{x}_i) - f(\b{x}_j)\|_2^2$, $\beta = (1-\epsilon)$, and $\mu = (p/d) \|\b{x}_i - \b{x}_j\|_2^2$. As $0 < \epsilon < 1$, we have $\beta<1$. By Lemma \ref{lemma_P_beta_and_exp}, we have:
\begin{align}
&\mathbb{P}(\|f(\b{x}_i) - f(\b{x}_j)\|_2^2 \leq (1-\epsilon) (p/d) \|\b{x}_i - \b{x}_j\|_2^2) \nonumber\\
&\overset{(\ref{equation_P_beta_and_exp_1})}{\leq} \exp\big(\frac{p}{2} (1 - (1-\epsilon) + \ln (1-\epsilon))\big) \nonumber\\
&= \exp\big(\frac{p}{2} (\epsilon + \ln (1-\epsilon))\big) \overset{(a)}{\leq} \exp\Big(\frac{p}{2} \big(\epsilon - (\epsilon + \frac{\epsilon^2}{2})\big)\Big) \nonumber\\
&= \exp(-\frac{p\, \epsilon^2}{4}) \overset{(b)}{\leq} \exp(-2 \ln(n)) = \frac{1}{n^2},
\end{align}
where $(a)$ is because:
\begin{align*}
\ln(1-x) \leq -(x + \frac{x^2}{2}), \quad \forall x \geq 0,
\end{align*}
and $(b)$ is because:
\begin{align*}
&p \overset{(\ref{equation_JL_lemma_p_lower_bound_2})}{\geq} 4(\frac{\epsilon^2}{2} - \frac{\epsilon^3}{3})^{-1} \ln(n) \\
&~~~~~~~~~~ \implies \exp(-\frac{p\, \epsilon^2}{4}) \leq \exp(\frac{-6 \epsilon^2}{3\epsilon^2 - 2\epsilon^3} \ln(n)) \\
&~~~~~~~~~~~~~~~~~~~\overset{(c)}{\leq} \exp(-2 \ln(n)),
\end{align*}
where $(c)$ is because:
\begin{align*}
3\epsilon^2 \geq 3\epsilon^2 - 2\epsilon^3 \implies \frac{-6 \epsilon^2}{3\epsilon^2 - 2\epsilon^3} \leq -2.
\end{align*}
\textbf{Step 2 of proof:}
Consider one side of error which is $\mathbb{P}(\|f(\b{x}_i) - f(\b{x}_j)\|_2^2 \geq (1+\epsilon) (p/d) \|\b{x}_i - \b{x}_j\|_2^2)$. Take $L = \|f(\b{x}_i) - f(\b{x}_j)\|_2^2$, $\beta = (1+\epsilon)$, and $\mu = (p/d) \|\b{x}_i - \b{x}_j\|_2^2$. As $0 < \epsilon < 1$, we have $\beta>1$. By Lemma \ref{lemma_P_beta_and_exp}, we have:
\begin{align}
&\mathbb{P}(\|f(\b{x}_i) - f(\b{x}_j)\|_2^2 \geq (1+\epsilon) (p/d) \|\b{x}_i - \b{x}_j\|_2^2) \nonumber\\
&\overset{(\ref{equation_P_beta_and_exp_2})}{\leq} \exp\big(\frac{p}{2} (1 - (1+\epsilon) + \ln (1+\epsilon))\big) \nonumber\\
&= \exp\big(\frac{p}{2} (-\epsilon + \ln (1+\epsilon))\big) \nonumber\\
&\overset{(a)}{\leq} \exp\Big(\frac{p}{2} \big(-\epsilon+\epsilon-\frac{\epsilon^2}{2} + \frac{\epsilon^3}{3} \big)\Big) \nonumber\\
&= \exp\Big(\frac{-p}{2} \big(\frac{\epsilon^2}{2} - \frac{\epsilon^3}{3} \big)\Big) \overset{(b)}{\leq} \exp(-2 \ln(n)) = \frac{1}{n^2},
\end{align}
where $(a)$ is because:
\begin{align*}
\ln(1+x) \leq x - \frac{x^2}{2} + \frac{x^3}{3}, \quad \forall x > 0,
\end{align*}
and $(b)$ is because:
\begin{align*}
&p \overset{(\ref{equation_JL_lemma_p_lower_bound_2})}{\geq} 4(\frac{\epsilon^2}{2} - \frac{\epsilon^3}{3})^{-1} \ln(n) \\
&~~~~~~~~~~ \implies -\frac{p}{2} \leq -2(\frac{\epsilon^2}{2} - \frac{\epsilon^3}{3})^{-1} \ln(n) \\
&~~~~~~~~~~ \implies -\frac{p}{2} (\frac{\epsilon^2}{2} - \frac{\epsilon^3}{3}) \leq -2 \ln(n) \\
&~~~~~~~~~~ \implies \exp\Big(\!-\frac{p}{2} (\frac{\epsilon^2}{2} - \frac{\epsilon^3}{3})\Big) \leq \exp(-2 \ln(n)).
\end{align*}
Hence, we showed in both steps 1 and 2 that the probability of every side of error is bounded by $1/n^2$. Hence, the probability of error in random projection is of the order $\mathcal{O}(1/n^2)$. 
\end{proof}

\begin{definition}[With High Probability]
When the probability of success in an algorithm goes to one by increasing the input information or a parameter or parameters of the algorithm to infinity, the algorithm is correct with high probability, denoted by w.h.p.
\end{definition}

% The JL Lemma, with probability of success of Eq. (\ref{equation_JL_lemma_correct_projection}), is correct w.h.p. because $\delta \rightarrow 0$ in Eq. (\ref{equation_prob_error_JL_lemma}) if $p \rightarrow \infty$ (although note that $p$ is bounded from above by $d$). 
Lemma \ref{lemma_prob_success_whp} showed that the probability of success for random projection is of order $\mathcal{O}(1/n^2)$. 
The JL Lemma, whose probability of success is of order $\mathcal{O}(1/n^2)$, is correct w.h.p. because $\mathbb{P}(\text{error}) \rightarrow 0$ if $n \rightarrow \infty$ by increasing the number of data points. In other words, the more data points we have, the more accurate the random projection is in preserving local distances after projection onto the random subspace. 
Eq. (\ref{equation_JL_lemma_correct_projection}) shows that random projection does not distort the local distances by more than a factor of $(1 \pm \epsilon)$. Therefore, with a good probability, we have:
\begin{align}
\|\b{U}^\top (\b{x}_i - \b{x}_j)\|_2^2 \approx \|\b{x}_i - \b{x}_j\|_2^2,\,\, \forall i,j \in \{1, \dots, n\}.
\end{align}

It is shown in \cite{larsen2017optimality} that for any function $f(.)$ which satisfies Eq. (\ref{equation_JL_lemma_correct_projection}), the Eq. (\ref{equation_JL_lemma_p_lower_bound_relaxed}) holds for the lower bound on the dimensionality of subspace and it does not require $f(.)$ to be necessarily random projection. This shows that JL lemma is optimal. 
Moreover, it is shown in \cite{bartal2011dimensionality} that any finite subset of Euclidean space can be embedded in a subspace with dimensionality $\mathcal{O}(\epsilon^{-2} \ln(n))$ while the distances of points are preserved with at most $(1+\epsilon)$ distortion. 

\section{Sparse Linear Random Projection}\label{section_sparse_linear_random_projection}

The JL lemma, i.e. Theorem \ref{theorem_JL_lemma}, dealt with $\ell_2$ norm. There exist some works which formulate random projection and its proofs with $\ell_1$ norm to have sparse random projection. 
In the following, we introduce the lemmas and theorems which justify sparse random projection. The proofs of these theories are omitted for brevity and they can be found in the cited papers. 

The following lemma shows that if dataset is $k$-sparse, meaning that every vector has at most $k$ non-zero values, random projection preserves distances well enough. As the dataset is sparse, random projection is sparse.
% \cite{ward2014dimension}, slide 18
\begin{lemma}[\cite{baraniuk2006johnson}]
Consider the subset of $k$-sparse signals defined as:
\begin{align}
\mathcal{S}_k := \{\b{x} \in \mathbb{R}^d : \#\{i : |x_i| > 0\} \leq k\},
\end{align}
where $\#\{.\}$ denotes the cardinality of set. If $p = \mathcal{O}(\epsilon^{-2} k \ln(d/k))$ and $f\!: \mathbb{R}^d \rightarrow \mathbb{R}^p$, $f\!: \b{x} \mapsto \b{U}^\top \b{x}$, then with high probability, we have:
\begin{equation}\label{equation_k_sparse_signal_correctness}
\begin{aligned}
(1 - \epsilon) \|\b{x}_i - \b{x}_j\|_2^2 \leq \|&f(\b{x}_i) - f(\b{x}_j)\|_2^2 \\
&\leq (1 + \epsilon) \|\b{x}_i - \b{x}_j\|_2^2,
\end{aligned}
\end{equation}
for all $\b{x}_i, \b{x}_j \in \mathcal{S}_k$. The probability of error from Eq. (\ref{equation_k_sparse_signal_correctness}) is $2(12/\epsilon)^k e^{-(p/2) (\epsilon^2/8 - \epsilon^3/24)}$. 
\end{lemma}
% \begin{proof}
% Refer to \cite{baraniuk2006johnson}.
% \end{proof}

In the following, we define the Restricted Isometry Property (RIP) and show that a random projection matrix whose elements are sampled from a Gaussian distribution has this property.
% \cite{ward2014dimension}, slide 22
\begin{definition}[Restricted Isometry Property \cite{candes2006stable}]
The mapping $f: \mathbb{R}^d \rightarrow \mathbb{R}^p$, $f: \b{x} \mapsto \b{U}^\top \b{x}$ has the Restricted Isometry Property (RIP) of order $k$ and level $\epsilon \in (0,1)$ if:
\begin{align}
(1 - \epsilon) \|\b{x}\|_2^2 \leq \|\b{U}^\top \b{x}\|_2^2 \leq (1 + \epsilon) \|\b{x}\|_2^2,
\end{align}
for all $k$-sparse $\b{x} \in \mathbb{R}^d$.
\end{definition}

% \cite{ward2014dimension}, slide 22
\begin{lemma}[\cite{candes2006stable}]
A Gaussian random matrix $\b{U} \in \mathbb{R}^{d \times p}$ has RIP if $p \geq \epsilon^{-2} k \ln(d/k)$.
\end{lemma}

% \cite{ward2014dimension}, slide 27
\begin{theorem}[\cite{krahmer2011new}]\label{theorem_sparse_diagonal_sign_flip}
Consider the linear mapping $f: \mathbb{R}^d \rightarrow \mathbb{R}^p$, $f: \b{x} \mapsto \b{U}^\top \b{x}$ satisfying: 
\begin{align}
(1 - \epsilon) \|\b{x}\|_2^2 \leq \|\b{U}^\top \b{x}\|_2^2 \leq (1 + \epsilon) \|\b{x}\|_2^2,
\end{align}
for all $k$-sparse $\b{x} \in \mathbb{R}^d$ and suppose $\b{D}_\xi \in \mathbb{R}^{d \times d}$ is a diagonal matrix with $\xi = \pm 1$ on its diagonal. We have:
\begin{align}
\mathbb{P}\Big((1-\epsilon) \|\b{x}\|_2^2 \leq \|\b{U}^\top &\b{D}_\xi \b{x}\|_2^2 \leq (1+\epsilon) \|\b{x}\|_2^2\Big) \nonumber\\
&\geq 1 - 2e^{-c \epsilon^2 p / \ln(d)}.
\end{align}
\end{theorem}

\begin{corollary}
Comparing Theorem \ref{theorem_sparse_diagonal_sign_flip} with Definition \ref{definition_concentration_of_measure} shows that a random matrix with RIP and random column sign flips implies the JL lemma up to an $\ln(d)$ factor. 
\end{corollary}

The following lemma is going to be used in the proof of Corollary \ref{corollary_l1_meaningful_than_l2}. That corollary shows that $\ell_1$ norm makes more sense as a distance than the Euclidean distance for high-dimensional data. 
\begin{lemma}[{\citep[Theorem 2]{hinneburg2000nearest}}]\label{lemma_distance_of_points_wrt_norms}
Suppose we have $n$ i.i.d. random points in $\mathbb{R}^d$. Let $\text{dist}_{\text{min},r}$ and $\text{dist}_{\text{max},r}$ denote the minimum and maximum distances of points from the origin, respectively, with respect to a metric with norm $\ell_r$. We have:
\begin{align}
\lim_{d \rightarrow \infty} \mathbb{E}[\text{dist}_{\text{max},r} - \text{dist}_{\text{min},r}] = d^{1/r - 1/2}.
\end{align}
\end{lemma}

\begin{corollary}[\cite{hinneburg2000nearest}]\label{corollary_l1_meaningful_than_l2}
For high dimensional data, $\ell_1$ norm (such as Manhattan distance) is more meaningful for nearest neighbor comparisons than $\ell_2$ norm (such as Euclidean distance). 
\end{corollary}
\begin{proof}
Lemma \ref{lemma_distance_of_points_wrt_norms} shows that for a metric with $\ell_2$ norm (i.e. Euclidean distance), we have $\lim_{d \rightarrow \infty} \mathbb{E}[\text{dist}_{\text{max},r} - \text{dist}_{\text{min},r}] = 1$; therefore, all points are equidistant from the origin up to a constant. For a metric by $\ell_r$ norm with $r > 2$, we have $1/r - 1/2 < 0$ so the limit gets very small for large $d$ values (when $d \rightarrow \infty$). Hence, all points are equidistant from origin completely. 
For $\ell_1$ norm, we have $\lim_{d \rightarrow \infty} \mathbb{E}[\text{dist}_{\text{max},r} - \text{dist}_{\text{min},r}] = d^{1/2}$; hence, relative difference of nearest and farthest points from origin increases with dimension. Q.E.D.
\end{proof}

The above corollary showed that $\ell_1$ norm is more meaningful for nearest neighbor comparisons than $\ell_2$ norm. The below theorem extend the JL lemma to any $\ell_r$ norm with $1 \leq r \leq \infty$. 
% \cite{ward2014dimension}, slide 37
\begin{theorem}[Extension of JL lemma to any norm {\citep[Lemma 3.1]{lee2005metric}}, \cite{ward2014dimension}]\label{theorem_linear_mapping_lr_norm}
For embedding into a subspace equipped with $\ell_r$ norm with $1 \leq r \leq \infty$, there exists a linear mapping $f: \mathbb{R}^d \rightarrow \mathbb{R}^p$, $f: \b{x} \mapsto \b{U}^\top \b{x}$, with $p \leq d$, which satisfies:
\begin{align}
\big(\frac{p}{d}\big)^{|1/r - 1/2|} \|\b{x}_i &- \b{x}_j\|_r \leq \|\b{U}^\top (\b{x}_i - \b{x}_j)\|_r \nonumber\\
&\leq \big(\frac{d}{p}\big)^{|1/r - 1/2|} \|\b{x}_i - \b{x}_j\|_r,
\end{align}
for all $\b{x}_i, \b{x}_j \in \mathbb{R}^d$.
\end{theorem}

One may ask why we often use the $\ell_2$ norm, and not $\ell_1$ norm, for random projection. This is because the following corollary shows that the suitable norm for random projection is $\ell_2$ norm and it is very hard to use other norms for random projection \cite{charikar2002dimension,brinkman2005impossibility}. 
% \cite{ward2014dimension}, slide 37
\begin{corollary}[{\citep[Section 3]{lee2005metric}}]\label{corollary_impossible_l1_embedding}
The suitable norm for linear random projection is the $\ell_2$ norm. We have a difficulty or impossibility of dimensionality reduction with a linear map in $\ell_r$ with $r \neq 2$.
\end{corollary}
\begin{proof}
According to Theorem \ref{theorem_linear_mapping_lr_norm}, the local distances of points after projection is completely sandwiched between their local distances in the input space before projection; hence, $\|\b{U}^\top (\b{x}_i - \b{x}_j)\|_2 \approx \|\b{x}_i - \b{x}_j\|_2$. This shows that embedding into $\ell_2$ norm works very well, with constant distortion, in terms of local distance preservation. Consider $\ell_1$ as an example for $\ell_r$ with $r \neq 2$. In this case, we have $(p/d)^{1/2} \|\b{x}_i - \b{x}_j\|_1 \leq \|\b{U}^\top (\b{x}_i - \b{x}_j)\|_1 \leq (p/d)^{1/2} \|\b{x}_i - \b{x}_j\|_1$ which is not a constant distortion. Q.E.D.
\end{proof}

% \cite{ward2014dimension}, slide 40
\begin{definition}[RIP-1 {\citep[Definition 8]{berinde2008combining}}]
For $p \geq c\, \epsilon^{-2} k \ln(d)$, there exists a linear mapping $f: \mathbb{R}^d \rightarrow \mathbb{R}^p$, $f: \b{x} \mapsto \b{U}^\top \b{x}$, with $p \leq d$, which satisfies:
\begin{align}
(1 - 2\epsilon) \|\b{x}\|_1 \leq \|\b{U}^\top \b{x}\|_1 \leq \|\b{x}\|_1,
\end{align}
for all $k$-sparse $\b{x} \in \mathbb{R}^d$, where $\epsilon$ is the error tolerance and $c$ is a constant. 
If this holds for a matrix $\b{U}$, it has the Restricted Isometry Property-1 (RIP-1).
\end{definition}

We saw that in $\ell_1$ norm, random projection is hard to work well and in $\ell_2$ norm, it is not sparse. Hence, we need some norm in between $\ell_1$ and $\ell_2$ norms to have sparse random projection which performs well enough and is not that hard to be performed. 
In the following, we define an interpolation norm. 
\begin{definition}[Interpolation Norm {\citep[Appendix]{krahmer2016unified}}]
Suppose we rearrange the elements of data $\b{x} \in \mathbb{R}^d$ from largest to smallest values. 
We divide the support of $\b{x}$ into $s$ disjoint subsets $\{\mathcal{S}_i\}_{i=1}^s$ where $\mathcal{S}_1$ corresponds to the largest values of $\b{x}$. Let $\b{x}_{\mathcal{S}_i}$ denote the values of $\b{x}$ which fall in $\mathcal{S}_i$ represented as a vector. 
The interpolation norm, denoted by $\ell_{1,2,s}$, is defined as:
\begin{align}
\|\b{x}\|_{1,2,s} := \sqrt{\sum_{i=1}^{\lceil n/s \rceil} \|\b{x}_{\mathcal{S}_i}\|_1^2}.
\end{align}
\end{definition}
Note that if $s=1$ and $s=d$, we have $\|\b{x}\|_{1,2,1} = \|\b{x}\|_2$ and $\|\b{x}\|_{1,2,d} = \|\b{x}\|_1$, respectively. Also, for any $s$, if $\b{x}$ is $k$-sparse, we have $\|\b{x}\|_{1,2,s} = \|\b{x}\|_1$. 

The following theorem is a version of JL lemma which makes use the interpolation norm. 
\begin{theorem}[{\citep[Theorem 2]{krahmer2016unified}}]\label{theorem_linear_mapping_interpolation_norm}
There exists the mapping $f: \mathbb{R}^d \rightarrow \mathbb{R}^p$, $f: \b{x} \mapsto \b{U}^\top \b{x}$ such that with probability $(1 - 2d e^{-\epsilon^2 p / s})$, we have:
\begin{align}
(0.63 - \epsilon) \|\b{x}\|_{1,2,s} \leq \|\b{U}^\top \b{x}\|_1 \leq (1.63 + \epsilon) \|\b{x}\|_{1,2,s}.
\end{align}
\end{theorem}
For $s=1$ ($\ell_2$ norm) and $s=d$ ($\ell_1$ norm), Theorem \ref{theorem_linear_mapping_interpolation_norm} is reduced to Definition \ref{definition_concentration_of_measure} (random projection) and Theorem \ref{theorem_linear_mapping_lr_norm} with $\ell_1$ norm, respectively. 
For $s$-sparse $\b{x}$, Theorem \ref{theorem_linear_mapping_interpolation_norm} is reduced to embedding into a subspace equipped with $\ell_1$ norm. 
According to \cite{krahmer2016unified}, in $s=d$ case ($\ell_1$ norm) in Theorem \ref{theorem_linear_mapping_interpolation_norm}, we have $p \geq c d \ln(d)$ where $c$ is a constant. This contradicts $p \leq d$ so this coincides with Corollary \ref{corollary_impossible_l1_embedding}. 
However, for the $s$-sparse data in Theorem \ref{theorem_linear_mapping_interpolation_norm}, we need $p = \mathcal{O}(s \ln(d)) < \mathcal{O}(d)$ \cite{krahmer2016unified}. This shows that we can embed data onto a subspace equipped with $\ell_1$ norm, using linear projection, if the data are sparse enough. 

It is shown in \cite{baraniuk2006johnson} that the sparse random projection and JL lemma are related to compressed sensing \cite{donoho2006compressed}.
Some other works on sparse random projection are sparse Cauchy random projections \cite{li2007nonlinear,ramirez2012reconstruction}, Bernoulli random projection \cite{achlioptas2003database}, and very sparse random projection \cite{li2006very}. 
The sparse random projection can also be related to random projection on hypercubes {\citep[Chapter 7.2]{vempala2005random}}. This will be explained in Section \ref{section_random_projection_hypercube}. 

\section{Applications of Linear Random Projection}\label{section_applications_of_random_projection}

There exists various applications for random projection. Two moet well-known applications are low-rank matrix approximation and approximate nearest neighbor search which we explain in the following. 

\subsection{Low-Rank Matrix Approximation Using Random Projection}

% https://en.wikipedia.org/wiki/Low-rank_approximation

One of the applications of random projection is low-rank approximation of a matrix which we explain in the following.
According to the Eckart-Young-Mirsky theorem \cite{eckart1936approximation}, Singular Value Decomposition (SVD) can be used for low-rank approximation of a matrix. Consider the matrix $\b{X} \in \mathbb{R}^{d \times n}$ where $d \geq n$. 
The time complexity for SVD of this matrix $\mathcal{O}(d n^2)$ which grows if the matrix is large. We can improve the time complexity to $\mathcal{O}(d n \ln(n))$ using random projections \cite{papadimitriou2000latent}. In this technique, the low-rank approximation is not completely optimal but good enough. 
This method has two steps: First, it finds a smaller matrix $\b{Y}$ by random projection:
\begin{align}\label{equation_lowRankApprox_Y_projection}
\mathbb{R}^{p \times n} \ni \b{Y} := \frac{1}{\sqrt{p}} \b{U}^\top \b{X},
\end{align}
where $\b{U} \in \mathbb{R}^{d \times p}$ is the random projection matrix whose elements are independently drawn from standard normal distribution. Note that $p$ should satisfy Eq. (\ref{equation_JL_lemma_p_lower_bound}); for example, $p \geq c \ln(n) / \epsilon^2$ where $c$ is a positive constant.
As we have $p \ll d$, the matrix $\b{Y}$ is much smaller than the matrix $\b{X}$. We calculate the SVD of $\b{Y}$:
\begin{align}\label{equation_lowRankApprox_Y}
\b{Y} = \b{A} \b{\Lambda} \b{B}^\top = \sum_{i=1}^p \lambda_i \b{a}_i \b{b}_i^\top,
\end{align}
where $\b{A} = [\b{a}_1, \dots, \b{a}_p] \in \mathbb{R}^{p \times p}$ and $\b{B} = [\b{b}_1, \dots, \b{b}_p] \in \mathbb{R}^{n \times p}$ are the matrices of singular vectors and $\b{\Lambda} = \textbf{diag}([\lambda_1, \dots, \lambda_p]^\top) \in \mathbb{R}^{p \times p}$ contains the singular values. 
Note that the SVD of $\b{Y}$ is much faster than the SVD of $\b{X}$ because of the smaller size of matrix. 
The matrix $\b{X}$ can be approximated as its projection as:
\begin{align}\label{equation_low_rank_approx_X}
\mathbb{R}^{d \times n} \ni \widetilde{\b{X}}_p \approx \b{X} \Big(\sum_{i=1}^p \b{b}_i \b{b}_i^\top\Big).
\end{align}
Note that the right singular values $\b{b}_i$'s, which are used here, are equal to the eigenvectors of $\b{Y}^\top \b{Y}$ \cite{ghojogh2019unsupervised}. 
The rank of Eq. (\ref{equation_low_rank_approx_X}) is $p$ because $p \ll d, n$; hence, it is a low-rank approximation of $\b{X}$. 
In the following, we show why this is a valid low-rank approximation.

\begin{lemma}[{\citep[Lemma 3]{papadimitriou2000latent}}]
Let the SVD of $\b{X}$ be $\b{X} = \b{C} \b{\Sigma} \b{E}^\top = \sum_{i=1}^{d} \sigma_i \b{c}_i \b{e}_i^\top$ where $\b{C} = [\b{c}_1, \dots, \b{c}_{d}] \in \mathbb{R}^{d \times d}$, $\b{E} = [\b{e}_1, \dots, \b{e}_{d}] \in \mathbb{R}^{n \times d}$, and $\b{\Sigma} = \textbf{diag}([\sigma_1, \dots, \sigma_{d}]^\top) \in \mathbb{R}^{d \times d}$ are the left singular vectors, right singular vectors, and singular values, respectively. 
Also, let the SVD of $\b{X}$ with top $p$ singular values be $\b{X}_p := \sum_{i=1}^{p} \sigma_i \b{c}_i \b{e}_i^\top$. 
If $p \geq c \ln(n) / \epsilon^2$, the singular values of $\b{Y}$ are not much smaller than the singular values of $\b{X}$, i.e.:
\begin{align}\label{equation_lowRankApprox_lambda_squared_geq}
\sum_{i=1}^p \lambda_i^2 \geq (1 - \epsilon) \sum_{i=1}^p \sigma_i^2 = (1 - \epsilon) \|\b{X}_p\|_F^2,
\end{align}
where $\|.\|_F$ is the Frobenius norm. 
\end{lemma}
\begin{proof}
Refer to {\citep[Appendix]{papadimitriou2000latent}}.
\end{proof}

\begin{theorem}[{\citep[Theorem 5]{papadimitriou2000latent}}]
The low-rank approximation $\widetilde{\b{X}}$ approximates $\b{X}$ well enough:
\begin{align}
\|\b{X} - \widetilde{\b{X}}_p\|_F^2 \leq \|\b{X} - \b{X}_p\|_F^2 + 2 \epsilon \|\b{X}_p\|_F^2.
\end{align}
\end{theorem}
\begin{proof}
Consider the complete SVD with $\b{A} = [\b{a}_1, \dots, \b{a}_p] \in \mathbb{R}^{p \times p}$, $\b{B} = [\b{b}_1, \dots, \b{b}_{n}] \in \mathbb{R}^{n \times n}$, and $\b{\Lambda} \in \mathbb{R}^{p \times n}$.
We have:
\begin{align}
&\|\b{X} - \widetilde{\b{X}}_p\|_F^2 \overset{(a)}{=} \sum_{i=1}^{n} \|(\b{X} - \widetilde{\b{X}}_p) \b{b}_i\|_2^2 \nonumber\\
&= \sum_{i=1}^{n} \|\b{X} \b{b}_i - \widetilde{\b{X}}_p \b{b}_i \|_2^2 
\nonumber\\
&\overset{(\ref{equation_low_rank_approx_X})}{=} \sum_{i=1}^{n} \|\b{X} \b{b}_i - \b{X} \big(\sum_{i=1}^p \b{b}_i \underbrace{\b{b}_i^\top\big) \b{b}_i}_{=1}\|_2^2 \overset{(b)}{=} \sum_{i=p+1}^{n} \|\b{X} \b{b}_i\|_2^2 \nonumber\\
&= \|\b{X}\|_F^2 - \sum_{i=1}^p \|\b{X} \b{b}_i\|_2^2, \label{equation_lowRankApprox_X_minus_Xtilde}
\end{align}
where $(a)$ and $(b)$ are because the singular vectors $\{\b{a}_i\}$ and $\{\b{b}_i\}$ are orthonormal. 
Similarly, we can show:
\begin{align}
\|\b{X} - \b{X}_p\|_F^2 = \|\b{X}\|_F^2 - \|\b{X}_p\|_F^2. \label{equation_lowRankApprox_X_minus_Xp}
\end{align}
From Eqs. (\ref{equation_lowRankApprox_X_minus_Xtilde}) and (\ref{equation_lowRankApprox_X_minus_Xp}), we have:
\begin{align}\label{equation_lowRankApprox_X_minus_Xp_tilde_equals_sth}
\|\b{X} - \widetilde{\b{X}}_p\|_F^2 = \|\b{X} - \b{X}_p\|_F^2 + \|\b{X}_p\|_F^2 - \sum_{i=1}^p \|\b{X} \b{b}_i\|_2^2.
\end{align}
We also have:
\begin{align}
&\sum_{i=1}^p \|\b{Y} \b{b}_i\|_2^2 = \sum_{i=1}^p \b{b}_i^\top \b{Y}^\top \b{Y} \b{b}_i \nonumber\\
&\overset{(\ref{equation_lowRankApprox_Y})}{=} \sum_{i=1}^p \b{b}_i^\top \sum_{j=1}^p \lambda_j \b{b}_j \b{a}_j^\top \sum_{k=1}^p \lambda_k \b{a}_k \b{b}_k^\top \b{b}_i \nonumber\\
&= \sum_{i=1}^p \sum_{j=1}^p \sum_{k=1}^p \lambda_j \lambda_k \underbrace{\b{b}_i^\top \b{b}_j}_{=\delta_{ij}} \underbrace{\b{a}_j^\top \b{a}_k}_{=\delta_{jk}} \underbrace{\b{b}_k^\top \b{b}_i}_{=\delta_{ki}} \overset{(a)}{=} \sum_{i=1}^p \lambda_i^2, \label{equation_lowRankApprox_sum_lambda_squared}
\end{align}
where $\delta_{ij}$ is the Kronecker delta and $(a)$ is because the singular vectors $\{\b{a}_i\}$ and $\{\b{b}_i\}$ are orthonormal. 
On the other hand, we have:
\begin{align}
&\sum_{i=1}^p \|\b{Y} \b{b}_i\|_2^2 \overset{(\ref{equation_lowRankApprox_Y_projection})}{=} \sum_{i=1}^p \|\frac{1}{\sqrt{p}} \b{U}^\top \b{X} \b{b}_i\|_2^2 \nonumber \\
&\overset{(\ref{equation_lowRankApprox_sum_lambda_squared})}{\implies} \lambda_i^2 = \sum_{i=1}^p \|\frac{1}{\sqrt{p}} \b{U}^\top \b{X} \b{b}_i\|_2^2. \label{equation_lowRankApprox_lambda_squared}
\end{align}
According to Eq. (\ref{equation_JL_lemma_correct_projection}), with high probability, we have:
\begin{align}\label{equation_lowRankApprox_JL_inequality}
(1-\epsilon) \|\b{X}\b{b}_i\|_2^2 \leq \|\frac{1}{\sqrt{p}} \b{U}^\top \b{X} \b{b}_i\|_2^2 \leq (1+\epsilon) \|\b{X}\b{b}_i\|_2^2.
\end{align}
From Eqs. (\ref{equation_lowRankApprox_lambda_squared}) and (\ref{equation_lowRankApprox_JL_inequality}), with high probability, we have:
\begin{align*}
&\lambda_i^2 \leq (1+\epsilon) \|\b{X}\b{b}_i\|_2^2 \\
&\implies \sum_{i=1}^p (1+\epsilon) \|\b{X}\b{b}_i\|_2^2 \geq \sum_{i=1}^p \lambda_i^2 \overset{(\ref{equation_lowRankApprox_lambda_squared_geq})}{\geq} (1 - \epsilon) \|\b{X}_p\|_F^2 \\
&\implies \sum_{i=1}^p \|\b{X}\b{b}_i\|_2^2 \geq \frac{1-\epsilon}{1+\epsilon} \|\b{X}_p\|_F^2 \geq (1-2\epsilon) \|\b{X}_p\|_F^2 \\
&\implies \|\b{X}_p\|_F^2 - \sum_{i=1}^p \|\b{X} \b{b}_i\|_2^2 \leq 2\epsilon\, \|\b{X}_p\|_F^2 \\
&\overset{(\ref{equation_lowRankApprox_X_minus_Xp_tilde_equals_sth})}{\implies} \|\b{X} - \widetilde{\b{X}}_p\|_F^2 \leq \|\b{X} - \b{X}_p\|_F^2 + 2\epsilon\, \|\b{X}_p\|_F^2.
\end{align*}
Q.E.D.
\end{proof}

\begin{lemma}[\cite{papadimitriou2000latent}]
The time complexity of low-rank approximation (with rank $p$) of matrix $\b{X}$ using random projection is $\mathcal{O}(d n \ln(n))$ which is much better than the complexity of SVD on $\b{X}$ which is $\mathcal{O}(d n^2)$. 
\end{lemma}
\begin{proof}
Complexities of Eqs. (\ref{equation_lowRankApprox_Y_projection}) and (\ref{equation_lowRankApprox_Y}) are $\mathcal{O}(d n p)$ and $\mathcal{O}(n p^2)$, respectively, noticing that $p \ll n \leq d$. We also have $p \geq c \ln(n) / \epsilon^2$ where $c$ is a constant. Hence, the total complexity is $\mathcal{O}(d n p) + \mathcal{O}(n p^2) \in \mathcal{O}(d n p) = \mathcal{O}(d n \ln(n))$. Q.E.D.
\end{proof}

% slides \cite{liberty2006random}, slide 24:
Some other works, such as \cite{frieze2004fast,martinsson2007fast}, have improved the time complexity of using random projection for low-rank applications.
For more information on using random projection for low-rank approximation, one can refer to {\citep[Chapter 8]{vempala2005random}}. 

\subsection{Approximate Nearest Neighbor Search}

\subsubsection{Random Projection onto Hypercube}\label{section_random_projection_hypercube}

Let $\mathbb{Z}^d_w$ denote the set of $d$-dimensional integer vectors with $w$ possible values; for example, $\mathbb{Z}^d_2 := \{0,1\}^d$ is the set of binary values. Also let $\|\b{x}_i - \b{x}_j\|_H$ denote the Hamming distance between two binary vectors $\b{x}_i$ and $\b{x}_j$. We can have random projection of a binary dataset onto a hypercube $\mathbb{Z}^d_w$ where $\ell_1$ norm or Hamming distances are used \cite{kushilevitz1998algorithm,kushilevitz2000efficient}. 

% slides \cite{liberty2006random}, slide 27
\begin{theorem}[Random projection onto hypercube \cite{kushilevitz1998algorithm,kushilevitz2000efficient}]\label{theorem_random_projection_hypercube}
Let $\text{mod}$ denote the modulo operation.
Consider a binary vector $\b{x} \in \mathbb{Z}^d_2$ which is projected as $f(\b{x}) = \b{U}^\top \b{x} \text{ mod } 2$ with a random binary projection matrix $\mathbb{R} \in \{0,1\}^{d \times p}$, where $p = O(\ln(n) / \epsilon^2)$ (usually $p \ll d$). The elements of $\b{U}$ are i.i.d. with Bernoulli distribution having probability $\xi = (\epsilon^2 / \ell)$ to be one and probability $(1-\xi)$ to be zero. We have:
\begin{align}
& \text{if }\,\, \|\b{x}_i - \b{x}_j\|_H < \frac{\ell}{4} \nonumber\\
&\quad\quad\quad\implies \|f(\b{x}_i) - f(\b{x}_j)\|_H < (1+\epsilon)p\,\xi\,\frac{\ell}{4}, \\
& \text{if }\,\, \frac{\ell}{4} \leq \|\b{x}_i - \b{x}_j\|_H \leq \frac{\ell}{2\epsilon} \nonumber\\
&\implies (1-\epsilon)p\,\xi\leq\frac{\|f(\b{x}_i) - f(\b{x}_j)\|_H}{\|\b{x}_i - \b{x}_j\|_H} < (1+\epsilon)p\,\xi, \\
& \text{if }\,\, \|\b{x}_i - \b{x}_j\|_H > \frac{\ell}{2\epsilon} \nonumber\\
&\quad\quad\quad\implies \|f(\b{x}_i) - f(\b{x}_j)\|_H > (1-\epsilon)p\,\xi\,\frac{\ell}{2\epsilon},
\end{align}
for all $\b{x}_i, \b{x}_j \in \mathbb{Z}^d_2$, with probability at least $(1-e^{-c \epsilon^4 p})$ where $c$ is a positive constant. 
\end{theorem}
\begin{proof}
Proof is available in \cite{vempala2005random,liberty2006random}.
\end{proof}

% \begin{proof}
% Proof is based on \cite{vempala2005random,liberty2006random}.
% We have:
% \begin{align*}
% &\mathbb{E}\big[\|f(\b{x}_i) - f(\b{x}_j)\|_H\big] = \mathbb{E}\big[\|\b{U}^\top (\b{x}_i - \b{x}_j) \text{ mod } 2\|_H\big] \\
% &= \mathbb{E}\Big[\sum_{t=1}^p \b{u}_t^\top (\b{x}_i - \b{x}_j) \text{ mod } 2\Big] \\
% &\overset{(a)}{=} p\, \mathbb{P}\Big[\b{u}_1^\top (\b{x}_i - \b{x}_j) = 1 \text{ mod } 2\Big].
% \end{align*}
% \end{proof}

\subsubsection{Approximate Nearest Neighbor Search by Random Projection}

Consider a dataset $\mathcal{X} := \{\b{x}_i \in \mathbb{R}^d\}_{i=1}^n$. 
The nearest neighbor search problem refers to finding the closest point of dataset $\b{x}^* \in \mathcal{X}$ to a query point $\b{q} \in \mathbb{R}^d$.
One solution to this problem is to calculate the distances of the query point from all $n$ points of dataset and return the point with the smallest distance. However, its time and space complexities are both $\mathcal{O}(nd)$ which are not good. There is an algorithm for nearest neighbor search \cite{meiser1993point} with time complexity $\mathcal{O}(\text{poly}(d, \ln(n)))$ where $\text{ploy}()$ is a polynomial combination of its inputs. However, the space complexity of this algorithm is $\mathcal{O}(n^d)$. 

To have better time and space complexities, approximate nearest neighbor search is used. Approximate nearest neighbor search returns a point $\b{x}^* \in \mathcal{X}$ which satisfies:
\begin{align}
\|\b{q} - \b{x}^*\|_2 \leq (1 + \epsilon)\, (\min_{x \in \mathcal{X}} \|\b{q} - \b{x}\|_2),
\end{align}
where $\epsilon >0$ is the error tolerance. This problem is named $\epsilon$-approximate nearest neighbor search problem \cite{andoni2006near}.
We can relax this definition if we take the acceptable distance $r$ from the user:
\begin{align}\label{equation_approx_knn_with_r}
\|\b{q} - \b{x}^*\|_2 \leq (1 + \epsilon)\, r.
\end{align}
If no such point is found in $\mathcal{X}$, null is returned. 
This relaxation is valid because the smallest $r$ can be found using a binary search whose number of iterations is constant with respect to $d$ and $n$. 

We can use the introduced random projection onto hypercube for $\epsilon$-approximate nearest neighbor search \cite{kushilevitz1998algorithm,kushilevitz2000efficient}. 
We briefly explain this algorithm in the following. 
Assume the dataset is binary, i.e., $\mathcal{X} := \{\b{x}_i \in \mathbb{Z}_2^d\}_{i=1}^n$. If not, the values of vector elements are quantized to binary strings and the values of a vector are reshaped to become a binary vector. 
Let the dimensionality of binary (or quantized and reshaped) vectors is $d$. 
Recall the binary search required for finding the smallest distance $r$ in Eq. (\ref{equation_approx_knn_with_r}). 
In the binary search for finding the smallest distance $r$, we should try several distances. For every distance, we perform $k$ independent random projections onto $k$ hypercubes and then we select one of these random projections randomly. In every random projection, the points $\mathcal{X}$ are projected onto a random $p$-dimensional hypercube to have $\{f(\b{x}_i) \in \mathbb{Z}_2^p\}_{i=1}^n$. 
In the low-dimensional projected subspace, the comparison of points is very fast because (I) the subspace dimensionality $p$ is much less than the original dimensionality $d$ and (II) calculation of Hamming distance is faster and easier than Euclidean distance. Therefore, random projections onto hypercubes are very useful for approximate nearest neighbor search. 

\begin{lemma}
The above algorithm for approximate nearest neighbor search is correct with probability $(1 - \delta)$ where $0< \delta \ll 1$. 
The time complexity of the above algorithm is $\mathcal{O}(\frac{d}{\epsilon^4} \ln(\frac{n}{\delta}) \ln(d))$ and its space complexity is $\mathcal{O}(d^2 (c_1 n \ln(d))^{c_2 / \epsilon^4})$ where $c_1$ and $c_2$ are constants. 
\end{lemma}
\begin{proof}
Proof is available in \cite{kushilevitz1998algorithm,kushilevitz2000efficient,vempala2005random,liberty2006random}.
\end{proof}

% \begin{proof}
% According to Theorem \ref{theorem_random_projection_hypercube}, in each of these projections, the probability of error is $e^{-c\epsilon^4 p}$ for every pair of points. With the same way we came up with Eq. (\ref{equation_prob_error_JL_lemma}) from Eq. (\ref{equation_JL_prob_error_twoPoints}) by the Bonferroni's inequality, we have the probability of error for the whole $n$ points as $\mathcal{O}(2ne^{-c\epsilon^4 p})$. Therefore, the expectation of the number of random projections which fail is $\mathbb{E}[\#\text{fails}] = $
% \end{proof}

Random projection for approximate nearest neighbor search is also related to hashing; for example, see locality sensitive hashing \cite{slaney2008locality}.
There exist some other works on random projection for approximate nearest neighbor search \cite{indyk1998approximate,ailon2006approximate}. 
For more information on using random projection for approximate nearest neighbor search, one can refer to {\citep[Chapter 7]{vempala2005random}}.

\section{Random Fourier Features and Random Kitchen Sinks for Nonlinear Random Projection}\label{section_RFF_and_RKS}

So far, we explained linear random projection in which a linear projection is used. We can have nonlinear random projection which is a much harder task to analyze theoretically. A nonlinear random projection can be modeled as a linear random projection followed by a nonlinear function. Two fundamental works on nonlinear random projection are RFF \cite{rahimi2007random} and RKS \cite{rahimi2008weighted}, explained in the following. 

\subsection{Random Fourier Features for Learning with Approximate Kernels}

When the pattern of data is nonlinear, wither a nonliner algorithm should be used or the nonlinear data should be transformed using kernels to be able to use the linear methods for nonlinear patterns \cite{ghojogh2021reproducing}. 
Computation of kernels is a time-consuming task because the points are pulled to the potentially high dimensional space and then the inner product of pulled points to the Reproducing Kernel Hilbert Space (RKHS) are calculated (see \cite{ghojogh2021reproducing} for more details). 
Random Fourier Features (RFF) are used for accelerating kernel methods \cite{rahimi2007random}. 
For this, RFF transforms data to a low-dimensional space in contrast to the kernels which transform data to a potentially high dimensional space. RFF approximates the kernel, which is the inner product of pulled data to RKHS, by inner product of low-dimensional feature maps of data. This feature map is a random feature map and is $z: \mathbb{R}^d \rightarrow \mathbb{R}^{2p}$ where $p \ll d$; therefore, its computation of inner products is much faster than computing the inner product of data in RKHS. It satisfies the following approximation:
\begin{align}\label{equation_RFF_kernel_approximation}
k(\b{x}, \b{y}) = \b{\phi}(\b{x})^\top \b{\phi}(\b{y}) \approx z(\b{x})^\top z(\b{y}),
\end{align}
where $k(.,.)$ denotes the kernel function and $\b{\phi}(.)$ is the pulling function to RKHS. 
Note that $z$ is a nonlinear mapping which can be seen as a nonlinear random projection. It is a linear random projection $f : \mathbb{R}^d \rightarrow \mathbb{R}^{p}$ followed by nonlinear sine and cosine functions. We will see the formulation of $z$ and $f$ functions later. 

RFF works with positive definite kernels which are shift-invariant (also called stationary kernels), i.e., $k(\b{x}, \b{y}) = k(\b{x} - \b{y})$ \cite{ghojogh2021reproducing}.
Consider the inverse Fourier transform of the kernel function $k$:
\begin{equation}\label{equation_kernel_inverse_Fourier_transform}
\begin{aligned}
k(\b{x} - \b{y}) = \int_{\mathbb{R}^d} \widehat{k}(\b{u}) e^{j \b{u}^\top (\b{x} - \b{y})} d\b{u} \\
\overset{(a)}{=} \mathbb{E}_{\b{u}}[\zeta_{\b{u}}(\b{x}) \zeta_{\b{u}}(\b{y})^*],
\end{aligned}
\end{equation}
where $\b{u} \in \mathbb{R}^d$ is the frequency, $j$ is the imaginary unit, the superscript $^*$ denotes conjugate transpose, $\mathbb{E}[.]$ is the expectation operator, $(a)$ is because we define $\mathbb{R} \ni \zeta_{\b{u}}(\b{x}) := e^{j \b{u}^\top \b{x}}$, and $\widehat{k}(\b{u})$ is the Fourier transform of kernel function:
\begin{align}\label{equation_kernel_Fourier_transform}
\widehat{k}(\b{u}) = \frac{1}{2\pi} \int_{\mathbb{R}^d} e^{-j \b{u}^\top \b{x}}\, k(\b{x})\, d\b{x}. 
\end{align}
In Eq. (\ref{equation_kernel_inverse_Fourier_transform}), the kernel function $k$ and the transformed kernel $\widehat{k}$ are real valued so the sine part of $e^{j \b{u}^\top (\b{x} - \b{y})}$ can be ignored (see Euler's equation) and can be replaced with $\cos(\b{u}^\top (\b{x} - \b{y}))$; hence, $\zeta_{\b{u}}(\b{x}) = \cos(\b{u}^\top \b{x})$. Let:
\begin{align}\label{equation_RFF_f_w}
\mathbb{R}^2 \ni z_u(\b{x}) := [\cos(\b{u}^\top \b{x}), \sin(\b{u}^\top \b{x})]^\top.
\end{align}
In Eq. (\ref{equation_kernel_inverse_Fourier_transform}), $\widehat{k}(\b{u})$ can be seen as a $d$-dimensional probability density function. 
We draw $p$ i.i.d. random projection vectors, $\{\b{u}_t \in \mathbb{R}^d\}_{t=1}^p$, from $\widehat{k}(\b{u})$. 
We define a normalized vector version of Eq. (\ref{equation_RFF_f_w}) as:
\begin{equation}\label{equation_RFF_z_vector}
\begin{aligned}
\mathbb{R}^{2p} \ni z(\b{x}) := \frac{1}{\sqrt{p}} [&\cos(\b{u}_1^\top \b{x}), \dots, \cos(\b{u}_p^\top \b{x}), \\
&\sin(\b{u}_1^\top \b{x}), \dots, \sin(\b{u}_p^\top \b{x})]^\top.
\end{aligned}
\end{equation}
If we take linear random projection as:
\begin{align}
\mathbb{R}^{p} \ni f(\b{x}) := \b{U}^\top \b{x} = [\b{u}_1^\top \b{x}, \dots, \b{u}_p^\top \b{x}]^\top,
\end{align}
with the projection matrix $\b{U} := [\b{u}_1, \dots, \b{u}_p] \in \mathbb{R}^{d \times p}$, we see that the function $z(.)$ is applying sine and cosine functions to a linear random projection $f(.)$ of data $\b{x}$.

According to $\cos(a-b) = \cos(a) \cos(b) + \sin(a) \sin(b)$, we have:
\begin{align}\label{equation_RFF_f_innser_products}
z(\b{x})^\top z(\b{y}) = \cos(\b{u}^\top (\b{x} - \b{y})).
\end{align}
Considering $\zeta_{\b{u}}(\b{x}) = \cos(\b{u}^\top \b{x})$, we have:
\begin{align}
k(\b{x}, \b{y}) &= k(\b{x} - \b{y}) \overset{(\ref{equation_kernel_inverse_Fourier_transform})}{=} \mathbb{E}_{\b{u}}[\zeta_{\b{u}}(\b{x}) \zeta_{\b{u}}(\b{y})^*] \nonumber \\
&= \mathbb{E}_{\b{u}}[z(\b{x})^\top z(\b{y})]. \label{equation_RFF_kernel_approx_expectation}
\end{align}
Thus, Eq. (\ref{equation_RFF_kernel_approximation}) holds in expectation. 
We show in the following that the variance of Eq. (\ref{equation_RFF_kernel_approx_expectation}) is small to have Eq. (\ref{equation_RFF_kernel_approximation}) as a valid approximation. 
According to Eq. (\ref{equation_RFF_f_innser_products}), we have $-1 \leq z(\b{x})^\top z(\b{y}) \leq 1$; hence, we can use the Hoeffding's inequality \cite{hoeffding1963probability}:
\begin{align}\label{equation_RFF_Hoeffding}
\mathbb{P}\big(|z(\b{x})^\top z(\b{y}) - k(\b{x}, \b{y})| \geq \epsilon\big) \leq 2 \exp(-\frac{p\,\epsilon^2}{2}),
\end{align}
which is a bound on the probability of error in Eq. (\ref{equation_RFF_kernel_approximation}) with error tolerance $\epsilon \geq 0$. 
Therefore, the approximation in Eq. (\ref{equation_RFF_kernel_approximation}) holds and we can approximate kernel computation with inner product in a random lower dimensional subspace. This improves the speed of kernel machine learning methods because $p$ is much less than both the dimensionality of input space, i.e. $d$, and the dimensionality of RKHS. 
Although the Eq. (\ref{equation_RFF_Hoeffding}) proves the correctness of approximation in RFF, the following theorem shows more strongly that the approximation in RFF is valid. 

\begin{theorem}[{\citep[Claim 1]{rahimi2007random}}]
Let $\mathcal{M}$ be a compact subset of $\mathbb{R}^d$ with diameter $\text{diam}(\mathcal{M})$. We have:
\begin{align}
&\mathbb{P}\Big(\sup_{\b{x}, \b{y} \in \mathcal{M}} |z(\b{x})^\top z(\b{y}) - k(\b{x}, \b{y})| \geq \epsilon\Big) \nonumber\\
&~~~~~~~~~\leq 2^8 \big(\frac{\sigma\, \text{diam}(\mathcal{M})}{\epsilon}\big)^2 \exp\big(\!-\frac{p\,\epsilon^2}{4(d+2)}\big),
\end{align}
\end{theorem}
where $\sigma^2 := \mathbb{E}_{\b{u}}[\b{u}^\top \b{u}]$ with probability density function $\widehat{k}(\b{u})$ used in expectation. The required lower bound on the dimensionality of random subspace is:
\begin{align}
p \geq \Omega\Big(\frac{d}{\epsilon^2} \ln\big(\frac{\sigma\, \text{diam}(\mathcal{M})}{\epsilon}\big)\Big).
\end{align}
\begin{proof}
Proof is available in {\citep[Appendix A]{rahimi2007random}}.
\end{proof}

In summary, the algorithm of RFF is as follows. Given a positive definite stationary kernel function $k$, we calculate its Fourier transform by Eq. (\ref{equation_kernel_Fourier_transform}). Then, we treat this Fourier transform as a probability density function and draw $p$ i.i.d. random projection vectors $\{\b{u}_t \in \mathbb{R}^d\}_{t=1}^p$ from that. For every point $\b{x}$, the mapping $z(\b{x})$ is calculated using Eq. (\ref{equation_RFF_z_vector}). Then, for every two points $\b{x}$ and $\b{y}$, their kernel is approximated by Eq. (\ref{equation_RFF_kernel_approximation}). The approximated kernel can be used in kernel machine learning algorithms. 
As was explained, RFF can be interpreted as nonlinear random projection because it applied nonlinear sine and cosine functions to linear random projections. 

Finally, it is noteworthy that there also exists another method for kernel approximation, named random binning features, which we refer the interested reader to \cite{rahimi2007random} for information about it. Similar to RFF, random binning features can also be interpreted as nonlinear random projection. 

It is shown in \cite{yang2012nystrom} that there is a connection between the Nystr{\"o}m method \cite{ghojogh2021reproducing} and RFF. 
Some other works on RFF are \cite{chitta2012efficient,sutherland2015error,sriperumbudur2015optimal,li2019towards}.

\subsection{Random Kitchen Sinks for Nonlinear Random Projection}

% https://people.eecs.berkeley.edu/~brecht/kitchensinks.html

Random Kitchen Sinks (RKS) \cite{rahimi2008weighted} was proposed after development of RFF \cite{rahimi2007random}. 
RKS is a nonlinear random projection where a nonlinear function is applied to a linear random projection. 
RKS models this nonlinear random projection as a random layer of neural network connecting $d$ neurons to $p$ neurons. 
A random network layer is actually a linear random projection followed by a nonlinear activation function:
\begin{align}\label{equation_RKS_layer}
g(\b{x}) := \phi\Big(\sum_{j=1}^d \sum_{t=1}^p u_{jt}\, x_j\Big) \overset{(\ref{equation_linear_random_projection})}{=} \phi(\b{U}^\top \b{x}) = \phi\big(f(\b{x})\big),
\end{align}
where $g(.)$ is the function representing the whole layer, $\phi(.)$ is the possibly nonlinear activation function, $f(.)$ is the linear random projection, $\b{U} \in \mathbb{R}^{d \times p}$ is the matrix of random weights of network layer, $u_{jt}$ is the $(j,t)$-th element of $\b{U}$, and $x_j$ is the $j$-th element of $\b{x} \in \mathbb{R}^d$. 
Therefore, a random network layer can be seen as a nonlinear random projection.
According to the universal approximation theorem \cite{hornik1989multilayer,huang2006universal}, the Eq. (\ref{equation_RKS_layer}) can fit any possible decision function mapping continuous inputs to a finite set of classes to any desired level of accuracy.  
According to the representer theorem \cite{aizerman1964theoretical}, we can state the function $g(.)$ in RKHS as \cite{ghojogh2021reproducing}:
\begin{align}\label{equation_RKS_representer_theorem}
g(\b{x}) = \sum_{i=1}^\infty \alpha_i\, \phi(\b{x}; w_i),
\end{align}
where $\{w_i\} \in \Omega$ are the parameters ($\Omega$ is the space of parameters), $\phi(.)$ is the pulling function to RKHS, $\{\alpha_i\}$ are the weights. 
Comparing Eqs. (\ref{equation_RKS_layer}) and (\ref{equation_RKS_representer_theorem}) shows that there is a connection between RKS and kernels \cite{rahimi2008uniform} if we consider the pulling function $\phi(.)$ as the activation function, the parameters $\{w_i\}$ as the parameters of activation function, and the weights $\{\alpha_i\}$ as the weights of network layer. 

Consider a classification task where $l_i$ is the label of $\b{x}_i$. The empirical risk between the labels and the output of network layer is:
\begin{align}\label{equation_RKS_empirical_risk}
R_e(g) := \frac{1}{n} \sum_{i=1}^n \ell\big(g(\b{x}_i), l_i\big),
\end{align}
where $\ell$ is a loss function. 
The true risk is:
\begin{align}\label{equation_RKS_true_risk}
R_t(g) := \mathbb{E}_{\{\b{x}_i, l_i\}}\Big[ \ell\big(g(\b{x}_i), l_i\big)\Big],
\end{align}
where $\mathbb{E}[.]$ is the expectation operator. 
We want to minimize the empirical risk; according to Eq. (\ref{equation_RKS_representer_theorem}), we have:
\begin{align}\label{equation_RKS_empirical_risk_minimize_twoVars}
\underset{\{w_i, \alpha_i\}_{i=1}^p}{\text{minimize}}\quad R_e\Big(\sum_{i=1}^p \alpha_i\, \phi(\b{x}; w_i)\Big),
\end{align}
which is a joint minimization over $\{w_i\}_{i=1}^p$ and $\{\alpha_i\}_{i=1}^p$. 
In RKS, we sample the parameters of activation function, $\{w_i\}_{i=1}^p$, randomly from a distribution $\mathbb{P}(w)$ in space $\Omega$. Then, we minimize the empirical risk over only the network weights $\{\alpha_i\}_{i=1}^p$. In other words, we eliminate the optimization variables $\{w_i\}_{i=1}^p$ by their random selection and minimize over only $\{\alpha_i\}_{i=1}^p$. In practice, the distribution $\mathbb{P}(w)$ can be any distribution in the space of parameters of activation functions. 
Let $\b{\alpha} := [\alpha_1, \dots, \alpha_p]^\top \in \mathbb{R}^p$. If $\{w_i\}_{i=1}^p$ are randomly sampled for all $i \in \{1, \dots, n\}$, suppose $\b{g}_i := [\phi(\b{x}_i, \b{w}_1), \dots, \phi(\b{x}_i, \b{w}_p)]^\top \in \mathbb{R}^p, \forall i$. 
According to Eqs. (\ref{equation_RKS_representer_theorem}), (\ref{equation_RKS_empirical_risk}), and (\ref{equation_RKS_empirical_risk_minimize_twoVars}), the optimization in RKS is:
\begin{equation}\label{equation_RKS_optimziation}
\begin{aligned}
& \underset{\b{\alpha} \in \mathbb{R}^p}{\text{minimize}}
& & \frac{1}{n} \sum_{i=1}^n \ell\big(\b{\alpha}^\top \b{g}_i, l_i\big) \\
& \text{subject to}
& & \|\b{\alpha}\|_\infty \leq \frac{c}{p},
\end{aligned}
\end{equation}
where $\|.\|_\infty$ is the maximum norm and $c$ is a positive constant. In practice, RKS relaxes the constraint of this optimization to a quadratic regularizer: 
\begin{equation}\label{equation_RKS_optimziation_practice}
\begin{aligned}
& \underset{\b{\alpha} \in \mathbb{R}^p}{\text{minimize}}
& & \frac{1}{n} \sum_{i=1}^n \ell\big(\b{\alpha}^\top \b{g}_i, l_i\big) + \lambda\|\b{\alpha}\|_2^2,
\end{aligned}
\end{equation}
where $\lambda >0$ is the regularization parameter. 
Henceforth, let the solution of Eq. (\ref{equation_RKS_optimziation}) or (\ref{equation_RKS_optimziation_practice}) be $\{\alpha_i\}_{i=1}^p$ and the randomly sampled parameters be $\{w_i\}_{i=1}^p$.
We put into into Eq. (\ref{equation_RKS_representer_theorem}) but with only $p$ components of summation to result the solution of RKS, denoted by $\widehat{g}'$:
\begin{align}\label{equation_RKS_final_solution}
\widehat{g}'(\b{x}) = \sum_{i=1}^p \alpha_i\, \phi(\b{x}; w_i).
\end{align}
In the following, we show that Eq. (\ref{equation_RKS_final_solution}) minimizes the empirical risk very well even when the parameters $\{w_i\}_{i=1}^p$ are randomly selected. Hence, RKS is a nonlinear random projection which estimates the labels $\{l_i\}_{i=1}^n$ with a good approximation. 

Consider the set of continuous-version of functions in Eq. (\ref{equation_RKS_representer_theorem}) as:
\begin{align}
\mathcal{G} := \Big\{g(\b{x}) = \int_\Omega \alpha(w)\, \phi(\b{x}; w)\, dw \, \Big| \, |\alpha(w)| \leq c\, \mathbb{P}(w) \Big\}.
\end{align}
Let $g^*$ be a function in $\mathcal{G}$ and $\{w_i\}_{i=1}^p$ be sampled i.i.d. from $\mathbb{P}(w)$. 
The output of RKS, which is Eq. (\ref{equation_RKS_final_solution}), lies in the random set:
\begin{align}
\widehat{\mathcal{G}} := \Big\{\widehat{g} = \sum_{i=1}^p \alpha_i\, \phi(\b{x}; w_i)\, \Big| \, |\alpha_i| \leq \frac{c}{p}, \forall i \Big\}.
\end{align}


% \begin{lemma}[{\citep[Lemma 1]{rahimi2008weighted}}]
% Let $\mu$ be a measure on $\mathbb{R}^d$ and $g^* \in \mathcal{G}$ be the minimizer of the true risk over $\mathcal{G}$. 
% For $\delta >0$, with probability at least $1-\delta$, there exists a function $\widehat{g} \in \widehat{\mathcal{G}}$ such that:
% \begin{align}
% \sqrt{\int_{\mathbb{R}^d} \big(\widehat{g}(\b{x}) - g^*(\b{x})\big)^2 d\mu(\b{x})} \leq \frac{c}{\sqrt{p}} \Bigg(1 + \sqrt{2 \ln(\frac{1}{\delta})}\Bigg).
% \end{align}
% \end{lemma}

\begin{lemma}[Bound on the approximation error {\citep[Lemma 2]{rahimi2008weighted}}]\label{lemma_RKS_approximation_error_bound}
Suppose the loss $\ell(l,l')$ is $L$-Lipschitz. Let $g^* \in \mathcal{G}$ be the minimizer of the true risk over $\mathcal{G}$.
For $\delta >0$, with probability at least $(1-\delta)$, there exists $\widehat{g} \in \widehat{G}$ satisfying:
\begin{align}
R_t(\widehat{g}) - R_t(g^*) \leq \frac{Lc}{\sqrt{p}} \Bigg(1 + \sqrt{2 \ln(\frac{1}{\delta}})\Bigg).
\end{align}
The term $R_t(\widehat{g}) - R_t(g^*)$ is the approximation error and the above equation is an upper-bound on it. 
\end{lemma}

\begin{lemma}[Bound on the estimation error {\citep[Lemma 3]{rahimi2008weighted}}]\label{lemma_RKS_estimation_error_bound}
Suppose the loss can be stated as $\ell(l,l') = \ell(l l')$ and it is $L$-Lipschitz. 
For $\delta >0$ and for all $\widehat{g} \in \widehat{\mathcal{G}}$, with probability at least $(1-\delta)$, we have:
\begin{align}\label{equation_RKS_estimation_error_bound}
|R_t(\widehat{g}) - R_e(\widehat{g})| \leq \frac{1}{\sqrt{n}} \Bigg(4Lc + 2|\ell(0)| + Lc \sqrt{\frac{1}{2} \ln(\frac{1}{\delta}})\Bigg).
\end{align}
The term $|R_t(\widehat{g}) - R_e(\widehat{g})|$ is the estimation error and the above equation is an upper-bound on it. 
\end{lemma}

\begin{theorem}[Bound on error of RKS {\citep[Theorem 1]{rahimi2008weighted}}]\label{theorem_RKS_error}
Let the activation function $\phi(.)$ be bounded, i.e., $\sup_{\b{x}, w}|\phi(\b{x}; w)| \leq 1$. 
Suppose the loss can be stated as $\ell(l,l') = \ell(l l')$ and it is $L$-Lipschitz. 
The output of RKS, which is Eq. (\ref{equation_RKS_final_solution}), is $\widehat{g}' \in \widehat{\mathcal{G}}$ satisfying:
\begin{align}
R_t(\widehat{g}') - \min_{g \in \mathcal{G}} R_t(g) \leq \mathcal{O}\Bigg(\Big(\frac{1}{\sqrt{n}} + \frac{1}{\sqrt{p}}\Big) Lc \sqrt{\ln(\frac{1}{\delta})}\Bigg),
\end{align}
with probability at least $(1 - 2\delta)$.
\end{theorem}
\begin{proof}
Let $g^* \in \mathcal{G}$ be a minimizer of true risk over $\mathcal{G}$, $\widehat{g}' \in \widehat{\mathcal{G}}$ be a minimizer of empirical risk over $\widehat{\mathcal{G}}$, and $\widehat{g}^* \in \widehat{\mathcal{G}}$ be a minimizer of true risk over $\widehat{\mathcal{G}}$.
We have:
\begin{align}
R_t(\widehat{g}')& - R_t(g^*) \overset{(a)}{=} R_t(\widehat{g}') - R_t(\widehat{g}^*) + R_t(\widehat{g}^*) - R_t(g^*) \nonumber\\
&\leq | R_t(\widehat{g}') - R_t(\widehat{g}^*) | + R_t(\widehat{g}^*) - R_t(g^*). \label{equation_RKS_risk_diff_0}
\end{align}
By Lemma \ref{lemma_RKS_estimation_error_bound}, with probability at least $(1-\delta)$, we have:
\begin{align}
& | R_t(\widehat{g}^*) - R_e(\widehat{g}^*) | \leq \epsilon_\text{est}, \label{equation_RKS_risk_diff_1} \\
& | R_t(\widehat{g}') - R_e(\widehat{g}') | \leq \epsilon_\text{est}, \label{equation_RKS_risk_diff_2}
\end{align}
where $\epsilon_\text{est}$ is the right-hand side of Eq. (\ref{equation_RKS_estimation_error_bound}):
\begin{align}\label{equation_RKS_epsilon_est}
\epsilon_\text{est} := \frac{1}{\sqrt{n}} \Bigg(4Lc + 2|\ell(0)| + Lc \sqrt{\frac{1}{2} \ln(\frac{1}{\delta}})\Bigg).
\end{align}
We said that $\widehat{g}'$ is the minimizer of empirical risk over $\widehat{\mathcal{G}}$ so:
\begin{align}\label{equation_RKS_risk_diff_3}
R_e(\widehat{g}') \leq R_e(\widehat{g}^*).
\end{align}
From Eqs. (\ref{equation_RKS_risk_diff_1}), (\ref{equation_RKS_risk_diff_2}), and (\ref{equation_RKS_risk_diff_3}), with probability at least $(1-\delta)$, we have:
\begin{align}
| R_t(\widehat{g}')& - R_t(\widehat{g}^*) | \leq 2 \epsilon_\text{est} \nonumber\\
&\overset{(\ref{equation_RKS_epsilon_est})}{=} \frac{2}{\sqrt{n}} \Bigg(4Lc + 2|\ell(0)| + Lc \sqrt{\frac{1}{2} \ln(\frac{1}{\delta}})\Bigg). \label{equation_RKS_risk_diff_4}
\end{align}
On the other hand, the second term in Eq. (\ref{equation_RKS_risk_diff_0}) is the approximation error and by Lemma \ref{lemma_RKS_approximation_error_bound}, with probability at least $(1-\delta)$, we have:
\begin{align}\label{equation_RKS_risk_diff_5}
R_t(\widehat{g}^*) - R_t(g^*) \leq \epsilon_\text{app} := \frac{Lc}{\sqrt{p}} \Bigg(1 + \sqrt{2 \ln(\frac{1}{\delta}})\Bigg).
\end{align}
Using Eqs. (\ref{equation_RKS_risk_diff_0}), (\ref{equation_RKS_risk_diff_4}), and (\ref{equation_RKS_risk_diff_5}) and by Bonferroni's inequality or the so-called union bound \cite{bonferroni1936teoria}, i.e. Eq. (\ref{equation_Bonferroni_inequality}), we have:
\begin{align*}
R_t(\widehat{g}')& - R_t(g^*) \leq 2\epsilon_\text{est} + \epsilon_\text{app} \\
&= \frac{2}{\sqrt{n}} \Bigg(4Lc + 2|\ell(0)| + Lc \sqrt{\frac{1}{2} \ln(\frac{1}{\delta}})\Bigg) \\
&~~~~~ + \frac{Lc}{\sqrt{p}} \Bigg(1 + \sqrt{2 \ln(\frac{1}{\delta}})\Bigg) \\
&= \mathcal{O}\Bigg(\Big(\frac{1}{\sqrt{n}} + \frac{1}{\sqrt{p}}\Big) Lc \sqrt{\ln(\frac{1}{\delta})}\Bigg),
\end{align*}
with probability at least $(1-\delta-
\delta) = (1-2\delta)$. Q.E.D.
\end{proof}



\section{Other Methods for Nonlinear Random Projection}\label{section_other_nonlinear_random_projections}

In addition to RFF and RKS, there exist other, but similar, approaches for nonlinear random projection, such as ELM, random weights in neural network, and ensemble of random projections. In the following, we introduce these approaches. 

\subsection{Extreme Learning Machine}

Extreme Learning Machine (ELM) was initially proposed for regression as a feed-forward neural network with one hidden layer \cite{huang2006extreme}. 
It was then improved for a multi-layer perceptron network \cite{tang2015extreme}. Its extension to multi-class classification was done in \cite{huang2011extreme}. 
ELM is a feed-forward neural network whose all layers except the last layer are random. Let the sample size of training data be $n$, their dimensionality be $d$-dimensional, the one-to-last layer have $d'$ neurons, and the last layer has $p$ neurons where $p$ is the dimensionality of labels. Note that in classification task, labels are one-hot encoded so $p$ is the number of classes. Every layer except the last layer has a possibly nonlinear activation function and behaves like an RKS; although, it is learned by backpropagation and not Eq. (\ref{equation_RKS_optimziation_practice}). 
Let $\b{\beta} \in \mathbb{R}^{d' \times p}$ be the matrix of weights for the last layer. If the $n$ outputs of one-to-last layer be stacked in the matrix $\b{H} \in \mathbb{R}^{d' \times n}$ and their target desired labels are $\b{T} \in \mathbb{R}^{p \times n}$, we desire to have $\b{\beta}^\top \b{H} = \b{T}$.
After randomly sampling the weights of all layers except $\b{\beta}$, ELM learns the weights $\b{\beta}$ by solving a least squares problem:
\begin{equation}\label{equation_ELM_optimization}
\begin{aligned}
& \underset{\b{\beta} \in \mathbb{R}^{d' \times p}}{\text{minimize}}
& & \|\b{\beta}^\top \b{H} - \b{T}\|_F^2 + \lambda\|\b{\beta}\|_F^2,
\end{aligned}
\end{equation}
where $\|.\|_F$ is the Frobenius norm and $\lambda>0$ is the regularization parameter. In other words, the last layer of ELM behaves like a (Ridge) linear regression \cite{hoerl1970ridge}. 

It is shown in \cite{huang2006universal} that the random weights in neural network with nonlinear activation functions are universal approximators \cite{hornik1989multilayer}. Therefore, ELM works well enough for any classification and regression task. 
This shows the connection between ELM and RKS because both work on random weights on network but with slightly different approaches. This connection can be seen more if we interpret ELM using kernels \cite{huang2014insight}. 
There exists several surveys on ELM, such as \cite{huang2011extreme2,huang2015trends}. 


\subsection{Randomly Weighted Neural Networks}

After proposal of RKS and ELM, it was shown that a feed-forward or convolutional neural network whose all layers are random also work very well \cite{jarrett2009best}. In other words, a stack of several RKS models works very well because a random network can be seen as a stack of several nonlinear random projections. As each of these nonlinear random layers has an upper-bound on their probability of error (see Theorem \ref{theorem_RKS_error}), the total probability of error in the network is also bounded. 
Note that it does not matter whether the network is feed-forward or convolutional because in both architectures, the output of activation functions are projection by the weights of layers. 
It is shown in \cite{saxe2011random} that the convolutional pooling architectures are frequency selective and translation invariant even if their weights are random; therefore, randomly weighed convolutional network work well enough. 
This also explains why random initialization of neural networks before backpropagation is a good and acceptable initialization. 
The randomly weighed neural network has been used for object recognition \cite{jarrett2009best,chung2016random} and face recognition \cite{cox2011beyond}. 
Some works have even put a step further and have made almost everything in network, including weights and hyper-parameters like architecture, learning rate, and number of neurons random \cite{pinto2009high}. 

\subsubsection{Distance Preservation by Deterministic Layers}

Consider a layer of neural network with $p$ neurons to $d$ neurons. Let $\b{U} \in \mathbb{R}^{d \times p}$ be the weight matrix of layer and $g(.)$ be the Rectified Linear Unit (ReLU) activation function \cite{glorot2011deep}. 
The layer can be modeled as a linear projection followed by ReLU activation function, i.e., $g(\b{U}^\top \b{x}_j)$.
The following lemma shows that for deterministic weights $\b{U}$, the distances are preserved. 
\begin{lemma}[Preservation of Euclidean distance by a layer \cite{bruna2013learning,bruna2014signal}]
Consider a network layer $\b{U} \in \mathbb{R}^{d \times p}$.
If the output of the previous layer for two points are $\b{x}_i, \b{x}_j \in \mathbb{R}^d$, we have:
\begin{align}
L_1 \|\b{x}_i - \b{x}_j\|_2 \leq \|g(\b{U}^\top \b{x}_i) -& g(\b{U}^\top \b{x}_j)\|_2 \nonumber\\
&\leq L_2 \|\b{x}_i - \b{x}_j\|_2, 
\end{align}
where $0 < L_1 \leq L_2$ are the Lipschitz constants. 
\end{lemma}
\begin{proof}
Proof is available in \cite{bruna2014signal}. 
\end{proof}

\subsubsection{Distance Preservation by Random Layers}

We showed that for deterministic weights, the distances are preserved. Here, we show that random weights also preserve the distances as well as the angles between points.  
Consider a network layer $\b{U} \in \mathbb{R}^{d \times p}$ whose elements are i.i.d. random values sampled from the Gaussian distribution. The activation function is ReLU and denoted by $g(.)$.
Suppose the output of the previous layer for two points are $\b{x}_i, \b{x}_j \in \mathbb{R}^d$. Assume the input data to the layer lie on a manifold $\mathcal{M}$ with the Gaussian mean width:
\begin{align*}
w_\mathcal{M} := \mathbb{E}[\sup_{\b{x}_i, \b{x}_j \in \mathcal{M}} \b{q}^\top (\b{x}_i - \b{x}_j)],
\end{align*}
where $\b{q} \in \mathbb{R}^d$ is a random vector whose elements are i.i.d. sampled from the Gaussian distribution. 
We assume that the manifold is normalized so the input data to the layer lie on a sphere. 
Let $\mathbb{B}_r^d \subset \mathbb{R}^d$ be the ball with radius $r$ in $d$-dimensional Euclidean space. Hence, we have $\mathcal{M} \subset \mathbb{B}_r^d$. 

\begin{theorem}[Preservation of Euclidean distance by a random layer {\citep[Theorem 3]{giryes2016deep}}]\label{theorem_random_layer_Euclidean_preservation}
Consider a random network layer $\b{U} \in \mathbb{R}^{d \times p}$ where the output of the previous layer for two points are $\b{x}_i, \b{x}_j \in \mathbb{R}^d$. Suppose $\mathcal{M} \subset \mathbb{B}_r^d$.
Assume the angle between $\b{x}_i$ and $\b{x}_j$ is denoted by:
\begin{align*}
\theta_{i,j} := \cos^{-1}\Big(\frac{\b{x}_i^\top \b{x}_j}{\|\b{x}_i\|_2 \|\b{x}_j\|_2}\Big),
\end{align*}
and satisfies $0 \leq \theta_{i,j} \leq \pi$. If $p \geq c\, \delta^{-4} w_\mathcal{M}$ with $c$ as a constant, with high probability, we have:
\begin{align}
\Big| &\|g(\b{U}^\top \b{x}_i) - g(\b{U}^\top \b{x}_j)\|_2^2 \nonumber\\
&- \Big( \frac{1}{2} \|\b{x}_i - \b{x}_j\|_2^2 + \|\b{x}_i\|_2 \|\b{x}_j\|_2 \psi(\b{x}_i, \b{x}_j) \Big) \Big| \leq \delta, \label{equation_random_layer_Euclidean_preservation}
\end{align}
where $\psi(\b{x}_i, \b{x}_j) \in [0,1]$ is defined as:
\begin{align}
\psi(\b{x}_i, \b{x}_j) := \frac{1}{\pi}\big(\!\sin(\theta_{i,j}) - \theta_{i,j} \cos(\theta_{i,j})\big).
\end{align}
\end{theorem}
\begin{proof}
Proof is available in {\citep[Appendix A]{giryes2016deep}}.
\end{proof}

\begin{theorem}[Preservation of angles by a random layer {\citep[Theorem 4]{giryes2016deep}}]\label{theorem_random_layer_angle_preservation}
Suppose the same assumptions as in Theorem \ref{theorem_random_layer_Euclidean_preservation} hold and let $\mathcal{M} \subset \mathbb{B}_1^d \setminus \mathbb{B}_\beta^d$ where $\delta \ll \beta^2 < 1$. 
Assume the angle between $g(\b{U}^\top \b{x}_i)$ and $g(\b{U}^\top \b{x}_j)$ is denoted by:
\begin{align*}
\theta'_{i,j} := \cos^{-1}\Big(\frac{(g(\b{U}^\top \b{x}_i))^\top g(\b{U}^\top \b{x}_j)}{\|g(\b{U}^\top \b{x}_i)\|_2 \|g(\b{U}^\top \b{x}_j)\|_2}\Big),
\end{align*}
and satisfies $0 \leq \theta'_{i,j} \leq \pi$. 
With high probability, we have:
\begin{align}\label{equation_random_layer_angle_preservation}
\Big| &\cos(\theta'_{i,j}) - \big(\!\cos(\theta_{i,j}) + \psi(\b{x}_i, \b{x}_j)\big) \Big| \leq \frac{15 \delta}{\beta^2 - 2 \delta}. 
\end{align}
\end{theorem}
\begin{proof}
Proof is available in {\citep[Appendix B]{giryes2016deep}}.
\end{proof}

\begin{corollary}[Preservation of distances and angles by a random layer {\citep[Corollary 5]{giryes2016deep}}]\label{corollary_random_layer_preservation_distances_corollary}
By a random layer with weights $\b{U}$ and ReLU activation function $g(.)$, for every two points $\b{x}_i$ and $\b{x}_j$ as inputs to the layer, we have with high probability that:
\begin{equation}\label{equation_random_layer_preservation_distances_corollary}
\begin{aligned}
\frac{1}{2} \|\b{x}_i - \b{x}_j\|_2^2 - \delta \leq\, &\|g(\b{U}^\top \b{x}_i) - g(\b{U}^\top \b{x}_j)\|_2^2 \\
&\leq \|\b{x}_i - \b{x}_j\|_2^2 + \delta.
\end{aligned}
\end{equation}
A similar expression can be stated for preserving the angle by the layer. 
\end{corollary}
\begin{proof}
% fig 2 in \cite{giryes2016deep}
The relation of $\psi(\b{x}_i, \b{x}_j)$ and $\cos(\b{x}_i, \b{x}_j)$ is depicted in Fig. \ref{figure_plot_randomNet} for $\theta_{i,j} \in [0, \pi]$. As this figure shows:
\begin{align*}
& \theta_{i,j}=0,\,\, \cos(\b{x}_i, \b{x}_j)=1 \implies \psi(\b{x}_i, \b{x}_j)=0, \\
& \theta_{i,j}=\pi,\,\, \cos(\b{x}_i, \b{x}_j)=0 \implies \psi(\b{x}_i, \b{x}_j)=1.
\end{align*}
Therefore, if two points are similar, i.e. their angle is zero, we have $\psi(\b{x}_i, \b{x}_j)=0$. Having $\psi(\b{x}_i, \b{x}_j)=0$ in Eqs. (\ref{equation_random_layer_Euclidean_preservation}) and (\ref{equation_random_layer_angle_preservation}) shows that when the two input points $\b{x}_i, \b{x}_j$ to a layer are similar, we almost have:
\begin{align*}
& \|g(\b{U}^\top \b{x}_i) - g(\b{U}^\top \b{x}_j)\|_2^2 \approx \|\b{x}_i - \b{x}_j\|_2^2, \\
& \cos(\theta'_{i,j}) \approx \cos(\theta_{i,j}). 
\end{align*}
This proves preserving both Euclidean distance and angle of points by a network layer with random weights and ReLU activation function. Q.E.D.
\end{proof}

\begin{figure}[!t]
\centering
\includegraphics[width=3.2in]{./images/plot_randomNet}
\caption{Relation of $\psi(\b{x}_i, \b{x}_j)$ with $\cos(\theta_{i,j})$ for $\theta_{i,j} \in [0, \pi]$.}
\label{figure_plot_randomNet}
\end{figure}

It is noteworthy that a survey on randomness in neural networks exists \cite{scardapane2017randomness}. 


\subsection{Ensemble of Random Projections}

There exist several ensemble methods for random projections where an ensemble of random projections is used. Ensemble methods benefit from model averaging \cite{hoeting1999bayesian} and bootstrap aggregating (bagging) \cite{breiman1996bagging}. Bagging reduces the estimation variance especially if the models are not correlated \cite{ghojogh2019theory}. As the random projection models are random, they are uncorrelated; hence, an ensemble of them can improve performance. 

The paper \cite{giryes2016deep}, introduced before, has stated that a similar analysis to the analysis of ReLU activation function can be done for max pooling \cite{scherer2010evaluation}.
Consider an ensemble of $m$ random projection matrices $\{\b{U}_j \in \mathbb{R}^{d \times p}\}_{j=1}^m$. We can transform data as a max pooling over this ensemble \cite{karimi2017ensembles}: 
\begin{align}
f(\b{x}) &:= \max\{f_1(\b{x}), \dots, f_m(\b{x})\} \nonumber\\
&= \max\{\b{U}_1^\top \b{x}, \dots, \b{U}_m^\top \b{x}\}.
\end{align}
This max pooling is performed element-wise, similar to element-wise operation of ReLU activation function introduced before. 

We can also have an ensemble of random projections for the classification task. One of the algorithms for this is \cite{schclar2009random} which we explain in the following. 
We sample projection matrices randomly where their elements are i.i.d. Also, we normalize the columns of projection matrices to have unit length.
We apply the ensemble of projection matrices $\{\b{U}_j \in \mathbb{R}^{d \times p}\}_{j=1}^m$ to the $n$ training data points $\b{X} \in \mathbb{R}^{d \times n}$ to have $\{f_1(\b{X}), \dots, f_m(\b{X})\} := \{\b{U}_1^\top \b{X}, \dots, \b{U}_m^\top \b{X}\}$. 
Let the labels of training data be denoted by $\{l_i\}_{i=1}^n$. 
We train an arbitrary classifier model using these projected data and their labels. In other words, the model is trained by $\{(\{\b{U}_1^\top \b{x}_i\}_{i=1}^n, \{l_i\}_{i=1}^n), \dots, (\{\b{U}_m^\top \b{x}_i\}_{i=1}^n, \{l_i\}_{i=1}^n)\}$. In the test phase, the test point is projected onto the column spaces of the $m$ projection matrices. These projections are fed to the classifier and the final prediction of label is found by majority voting. 
Note that there is another more sophisticated ensemble method for random projections \cite{cannings2015random} which divides dataset into disjoint sets. In each set an ensemble of random projections is performed. In each set, the validation error of random projections are calculated and the random projection with smallest validation error is selected. After training, majority voting is performed for predicting the test label. 

\section{Conclusion}\label{section_conclusion}

In this paper, we introduced the theory of linear and nonlinear random projections. We explained the JL lemma and its proof. Sparse random projection using $\ell_1$ norm, low-rank matrix approximation by random projection, and approximate nearest neighbor search by random projection onto hypercube were covered. Then, RFF, RKS, ELM, random neural networks, and ensemble of random projections were explained. 
For brevity, we did not some methods for random projection, such as bilateral random projection \cite{zhou2012bilateral} and dual random projection \cite{zhang2013recovering}. 

\section*{Acknowledgement}

The authors hugely thank Rachel Ward and Edo Liberty whose slides \cite{ward2014dimension,liberty2006random} partially covered some materials in this tutorial paper.  

% the instructors of ... course at the ... University (you can see their YouTube channel) whose lectures partly covered some materials mentioned in this tutorial paper. 

% \appendix



\bibliography{References}
\bibliographystyle{icml2016}



















\end{document}


% This document was modified from the file originally made available by
% Pat Langley and Andrea Danyluk for ICML-2K. This version was
% created by Lise Getoor and Tobias Scheffer, it was slightly modified
% from the 2010 version by Thorsten Joachims & Johannes Fuernkranz,
% slightly modified from the 2009 version by Kiri Wagstaff and
% Sam Roweis's 2008 version, which is slightly modified from
% Prasad Tadepalli's 2007 version which is a lightly
% changed version of the previous year's version by Andrew Moore,
% which was in turn edited from those of Kristian Kersting and
% Codrina Lauth. Alex Smola contributed to the algorithmic style files.
